% =============================================================================
% anexo_metodologico.tex — Anexo metodológico del Proyecto Provincias
%
% COMPILAR:  cd 04_Docs && latexmk -pdf anexo_metodologico.tex
%
% ANTES DE COMPILAR, regenerar las tablas y las cifras:
%   Rscript 02_Code/run_all.R
%   Rscript 02_Code/99_checks/auditoria_end_to_end.R
%   Rscript 02_Code/99_checks/diccionario_indicadores.R
%   Rscript 02_Code/99_checks/generar_tablas_anexo.R
%
% Ninguna cifra de este documento está escrita a mano. Todas entran por
% \input{tablas/...} o por las macros de tablas/cifras.tex, que produce el
% último script a partir de las salidas reales del pipeline.
% =============================================================================

\documentclass[11pt,a4paper]{article}
% =============================================================================
% preambulo.tex — Formato del anexo metodológico
%
% Deliberadamente sobrio: es un documento de referencia técnica, para consultar
% y para auditar, no una pieza de divulgación. Los colores son los del sistema
% visual del proyecto (02_Code/04_figuras/DISENO.md) para que el anexo y las
% figuras del informe se lean como una sola familia.
% =============================================================================

\usepackage[utf8]{inputenc}
\usepackage[T1]{fontenc}
\usepackage[spanish,es-noquoting,es-tabla]{babel}
\usepackage{lmodern}
\usepackage{microtype}

\usepackage[a4paper,margin=2.5cm,headheight=14pt]{geometry}

% Una ruta larga sin espacios no se puede partir y TeX prefiere desbordar la
% caja antes que estirar la línea. Con este margen de emergencia estira.
\emergencystretch=3em
\usepackage{booktabs}
\usepackage{tabularx}
\usepackage{longtable}
\usepackage{array}
\usepackage{ragged2e}
\usepackage{enumitem}
\usepackage{fancyhdr}
\usepackage{titlesec}
\usepackage{xcolor}
\usepackage{graphicx}
\usepackage[most]{tcolorbox}
\usepackage{caption}
\usepackage[hidelinks]{hyperref}

% --- Colores del sistema visual (DISENO.md §2) -------------------------------
\definecolor{azulmarino}{HTML}{4458A6}
\definecolor{baya}{HTML}{C24A49}
\definecolor{verdeselva}{HTML}{159449}
\definecolor{tintauno}{HTML}{292E38}
\definecolor{tintados}{HTML}{595E66}
\definecolor{tintatres}{HTML}{888C94}
\definecolor{grilla}{HTML}{DCDEE1}
\definecolor{zafreclaro}{HTML}{A1B4DB}

\hypersetup{colorlinks=true, linkcolor=azulmarino, urlcolor=azulmarino,
            citecolor=azulmarino}

% --- Títulos ------------------------------------------------------------------
\titleformat{\section}{\normalfont\Large\bfseries\color{tintauno}}{\thesection}{1em}{}
\titleformat{\subsection}{\normalfont\large\bfseries\color{tintauno}}{\thesubsection}{1em}{}
\titleformat{\subsubsection}{\normalfont\normalsize\bfseries\color{tintados}}{\thesubsubsection}{1em}{}
\titlespacing*{\section}{0pt}{2.2ex plus 1ex minus .2ex}{1.2ex plus .2ex}

% --- Encabezado y pie ---------------------------------------------------------
\pagestyle{fancy}
\fancyhf{}
\renewcommand{\headrulewidth}{0.4pt}
\fancyhead[L]{\footnotesize\color{tintatres} Proyecto Provincias · Anexo metodológico}
\fancyhead[R]{\footnotesize\color{tintatres}\nouppercase{\leftmark}}
\fancyfoot[C]{\footnotesize\color{tintatres}\thepage}

% --- Tablas -------------------------------------------------------------------
\renewcommand{\arraystretch}{1.15}
\captionsetup{font=small, labelfont={bf,color=tintados},
              textfont={color=tintados}, justification=raggedright,
              singlelinecheck=false}
\newcolumntype{Y}{>{\RaggedRight\arraybackslash}X}

% --- Cajas --------------------------------------------------------------------
% Tres roles, y solo tres: una regla que el pipeline aplica, un hallazgo de la
% auditoría y una advertencia sobre el dato. Más tipos de caja y el lector deja
% de distinguirlas.

\newtcolorbox{regla}[1][]{
  enhanced, breakable, colback=white, colframe=azulmarino,
  boxrule=0pt, leftrule=2.5pt, arc=0pt, outer arc=0pt,
  left=8pt, right=8pt, top=6pt, bottom=6pt,
  colbacktitle=white, fonttitle=\bfseries\small, coltitle=azulmarino,
  title={Regla del pipeline}, #1}

\newtcolorbox{hallazgo}[1][]{
  enhanced, breakable, colback=white, colframe=baya,
  boxrule=0pt, leftrule=2.5pt, arc=0pt, outer arc=0pt,
  left=8pt, right=8pt, top=6pt, bottom=6pt,
  colbacktitle=white, fonttitle=\bfseries\small, coltitle=baya,
  title={Hallazgo de la auditoría}, #1}

\newtcolorbox{advertencia}[1][]{
  enhanced, breakable, colback=grilla!30, colframe=grilla,
  boxrule=0.4pt, arc=1pt,
  left=8pt, right=8pt, top=6pt, bottom=6pt,
  colbacktitle=grilla!30, fonttitle=\bfseries\small, coltitle=tintados,
  title={Advertencia sobre el dato}, #1}

% --- Diagramas ----------------------------------------------------------------
% Los estilos del DAG se declaran aquí una sola vez. Cada tipo de nodo
% corresponde a una etapa del flujo y cada tipo de arista a una forma de
% dependencia; los .tex generados solo referencian estos nombres, nunca colores
% ni tamaños sueltos. La leyenda del documento (§2.2) describe exactamente esta
% tabla de estilos: si aquí cambia un color, allí cambia la leyenda.

\usepackage{tikz}
\usetikzlibrary{arrows.meta, positioning, calc, fit, backgrounds, shapes.geometric}

\tikzset{
  % --- Nodos, uno por etapa del flujo ---
  nodobase/.style={
    draw, rounded corners=1.5pt, line width=0.5pt,
    font=\scriptsize, inner sep=2.6pt, align=left,
    minimum height=4.6mm, anchor=west},
  nodofuente/.style={nodobase, draw=tintatres, fill=grilla!35, text=tintauno},
  nodoscript/.style={nodobase, draw=azulmarino, fill=azulmarino!8, text=tintauno},
  nododerivado/.style={nodobase, draw=zafreclaro, fill=zafreclaro!25, text=tintauno},
  nodorecurso/.style={nodobase, draw=tintatres, fill=white, text=tintauno},
  nodoseccion/.style={nodobase, draw=azulmarino, fill=azulmarino!15, text=tintauno},
  nodosalida/.style={nodobase, draw=verdeselva, fill=verdeselva!10, text=tintauno},
  % --- Aristas ---
  flujo/.style={-{Stealth[length=1.4mm, width=1.1mm]},
                draw=tintatres, line width=0.28pt, opacity=0.75},
  flujoexcepcion/.style={-{Stealth[length=1.4mm, width=1.1mm]},
                draw=baya, line width=0.34pt, densely dashed, opacity=0.9},
  % --- Etapas del diagrama de arquitectura ---
  % `text width` fijo: sin él, una etiqueta larga ensancha la caja hasta
  % solaparse con la vecina y las flechas quedan sin sitio donde dibujarse.
  cajabase/.style={rounded corners=2pt, line width=0.6pt, align=center,
                font=\footnotesize, text width=2.9cm, inner sep=4pt,
                minimum height=1.15cm, text=tintauno},
  etapa/.style={cajabase, draw=azulmarino, fill=azulmarino!8},
  almacen/.style={cajabase, draw=tintatres, fill=grilla!40},
  salida/.style={cajabase, draw=verdeselva, fill=verdeselva!10},
  paso/.style={-{Stealth[length=2mm, width=1.6mm]}, draw=azulmarino,
               line width=0.8pt},
  control/.style={-{Stealth[length=2mm, width=1.6mm]}, draw=baya,
               line width=0.7pt, densely dashed},
}

% --- Atajos -------------------------------------------------------------------
\newcommand{\ruta}[1]{\texttt{\small #1}}
\newcommand{\fn}[1]{\texttt{\small #1}}
% Muestra de estilo para la leyenda de §2.2: dibuja el propio nodo o la propia
% arista, para que la leyenda no pueda desincronizarse de lo que declara.
\newcommand{\muestranodo}[1]{%
  \raisebox{-0.6mm}{\begin{tikzpicture}[baseline]
    \node[#1, minimum width=9mm, minimum height=3.4mm] at (0,0) {};
  \end{tikzpicture}}}
\newcommand{\muestraarista}[1]{%
  \raisebox{0.4mm}{\begin{tikzpicture}[baseline]
    \draw[#1] (0,0) -- (0.9,0);
  \end{tikzpicture}}}

% Generado por 02_Code/99_checks/generar_tablas_anexo.R — no editar.
\newcommand{\cifraMunicipios}{88}
\newcommand{\cifraProvincias}{11}
\newcommand{\cifraInsumos}{45}
\newcommand{\cifraInsumosPresentes}{43}
\newcommand{\cifraInsumosAusentes}{2}
\newcommand{\cifraCurados}{22}
\newcommand{\cifraDerivados}{15}
\newcommand{\cifraAgregados}{7.730}
\newcommand{\cifraBanderasRojas}{4}
\newcommand{\cifraIndicadores}{330}
\newcommand{\cifraConFicha}{12}
\newcommand{\cifraHojasConVacios}{46}
\newcommand{\cifraFiguras}{814}
\newcommand{\cifraMapas}{110}
\newcommand{\cifraTablasDerivadas}{15}
\newcommand{\cifraColumnasDerivadas}{464}
\newcommand{\cifraSinLlave}{2}
\newcommand{\cifraProsaAuditada}{3.391}
\newcommand{\cifraProsaSinOrigen}{0}
\newcommand{\cifraProsaPct}{100,00}
\newcommand{\cifraPoderDeteccion}{95,3}
\newcommand{\cifraFalsasProbadas}{300}
\newcommand{\cifraPanelIndicadores}{39}


\title{\vspace{-1.5cm}\bfseries\color{tintauno}
  Anexo metodológico\\[2mm]
  \large\mdseries\color{tintados}
  Tratamiento de datos, fuentes, protocolos, limpieza y actualización}
\author{\normalsize\color{tintados} Proyecto Provincias · Planes Estratégicos
  Provinciales de Antioquia}
\date{\normalsize\color{tintatres}\today}

\begin{document}
\maketitle
\thispagestyle{fancy}

\begin{abstract}
\noindent
Este anexo documenta cómo se construye cada cifra que publican los informes
provinciales: de qué archivo sale, qué transformaciones sufre, con qué fórmula
se agrega y cómo se actualiza. Está escrito para tres lectores distintos: quien
tenga que \emph{defender} una cifra ante la Gobernación o ante un municipio,
quien tenga que \emph{actualizar} el diagnóstico dentro de un año, y quien tenga
que \emph{auditarlo} sin conocer el proyecto.

El principio que ordena todo el documento es uno solo: \textbf{ninguna cifra
publicada puede existir sin un archivo fuente y una fórmula declarada}. Las
secciones \ref{sec:trazabilidad} y \ref{sec:limitaciones} verifican
mecánicamente que se cumpla, y dicen dónde todavía no se cumple.
\end{abstract}

\tableofcontents
\newpage

% =============================================================================
\section{Alcance y arquitectura}
% =============================================================================

\subsection{Qué produce el flujo}

El diagnóstico territorial cubre \cifraProvincias{} provincias que agrupan
\cifraMunicipios{} de los 125 municipios de Antioquia, en diez secciones
temáticas por provincia. Para cada provincia el flujo produce diez libros de
Excel —uno por sección, con todas sus tablas— y \cifraFiguras{} archivos de
figura entre las \cifraProvincias{} provincias, de los cuales \cifraMapas{} son
mapas. Cada figura se exporta dos veces: en PNG a 300\,dpi para Word y en PDF
vectorial para imprenta.

\begin{regla}[title={Regla 0 — el dato manda sobre el informe}]
Cuando la fuente actual y el informe publicado difieren, el flujo entrega la
cifra de la fuente actual y lo declara. El informe de 2026 no es la referencia
de verdad: es una fotografía de las fuentes que existían al escribirlo.
\end{regla}

\subsection{El universo territorial}

La pertenencia de un municipio a una provincia y a una subregión sale siempre
del crosswalk territorial, nunca de la columna de texto que traiga cada fuente.
Varias fuentes traen su propia columna \fn{prov} o \fn{subreg}, y en al menos
dos casos verificados vienen desalineadas respecto del listado oficial.

\begin{table}[htbp]\centering
\caption{Las once provincias y su composición municipal}
\small
% Generado por 02_Code/99_checks/generar_tablas_anexo.R — no editar a mano.
\begin{tabularx}{\linewidth}{r>{\raggedright\arraybackslash}p{4.6cm}rX}
\toprule
\textbf{Id} & \textbf{Provincia} & \textbf{Mpios.} & \textbf{Subregiones que la componen} \\
\midrule
1 & Agroindustrial del Occidente & 6 & OCCIDENTE \\
2 & Bioenergética del Norte de Antioquia & 13 & NORDESTE, NORTE \\
3 & Del Río Grande & 5 & NORTE \\
4 & Turística y Agroecológica & 7 & OCCIDENTE \\
5 & Agua, Bosques y Turismo & 12 & ORIENTE \\
6 & Cartama & 11 & SUROESTE \\
7 & De la Paz & 5 & ORIENTE \\
8 & San Juan & 6 & SUROESTE \\
9 & Minero Agroecológica & 5 & NORDESTE \\
10 & Penderisco y Sinifana & 8 & OCCIDENTE, SUROESTE \\
11 & Área Metropolitana & 10 & VALLE DE ABURRA \\
\bottomrule
\end{tabularx}

\end{table}

\begin{advertencia}[title={Una errata que hay que conservar}]
El nombre de la provincia~5 aparece en varias fuentes como
«POVINCIA DEL AGUA, BOSQUES Y TURISMO», sin la \emph{r}. El listado maestro ya
se corrigió, pero otros archivos conservan la errata. El emparejamiento del
flujo normaliza el texto y acepta las dos grafías; corregir la constante sin
corregir todos los archivos dejaría a la provincia~5 sin municipios.
\end{advertencia}

% =============================================================================
\section{Convenciones del documento}
\label{sec:convenciones}
% =============================================================================

Un anexo de estructura de datos se lee salteado, buscando un dato concreto. Para
que eso funcione, la notación tiene que ser previsible. Todo lo que sigue está
declarado una sola vez, aquí, y aplicado en todo el documento.

\subsection{Notación tipográfica}

\noindent\begin{tabularx}{\linewidth}{l >{\RaggedRight\arraybackslash}X}
\toprule
\textbf{Forma} & \textbf{Qué señala} \\
\midrule
\ruta{01\_Data/00\_Inputs} & Ruta dentro del repositorio, siempre relativa a la raíz \\
\fn{leer\_derivado()} & Función, variable o nombre de columna del código \\
\textit{cursiva} & Término técnico en su primera aparición, o énfasis de contraste \\
\textbf{negrita} & Cifra o afirmación que el lector debería poder citar \\
«comillas angulares» & Cita literal de una etiqueta que aparece en una tabla publicada \\
\bottomrule
\end{tabularx}

\vspace{6pt}
Las cajas marcan tres cosas distintas y solo tres:

\vspace{4pt}
\noindent\begin{tabularx}{\linewidth}{l >{\RaggedRight\arraybackslash}X}
\toprule
\textbf{Caja} & \textbf{Qué contiene} \\
\midrule
Regla del pipeline & Una decisión que el código aplica siempre. Cambiarla exige cambiar el código, no el criterio de quien escribe. \\
Hallazgo de la auditoría & Un resultado de la verificación automática, con su cifra. Nada entra aquí sin CSV de respaldo. \\
Advertencia sobre el dato & Un límite de la fuente que el informe debe declarar al publicar la cifra. \\
\bottomrule
\end{tabularx}

\subsection{Estilos del diagrama}
\label{sec:estilos-dag}

Los DAG de la sección~\ref{sec:dag} se dibujan con estos estilos, declarados en
\ruta{04\_Docs/preambulo.tex}. Cada muestra de esta leyenda se dibuja con el
mismo estilo que usa el diagrama: la leyenda no puede quedar desincronizada de
lo que describe.

\vspace{4pt}
\noindent\begin{tabularx}{\linewidth}{c l >{\RaggedRight\arraybackslash}X}
\toprule
\textbf{Muestra} & \textbf{Estilo} & \textbf{Qué representa} \\
\midrule
\muestranodo{nodofuente}   & \fn{nodofuente}   & Archivo de entrada, tal como lo entrega la fuente \\
\muestranodo{nodoscript}   & \fn{nodoscript}   & Script del pipeline: una transformación con autor y código \\
\muestranodo{nododerivado} & \fn{nododerivado} & Derivado en Parquet, regenerable \\
\muestranodo{nodorecurso}  & \fn{nodorecurso}  & Recurso que una sección consume, sea derivado, curado o crudo \\
\muestranodo{nodoseccion}  & \fn{nodoseccion}  & Sección temática del diagnóstico \\
\muestranodo{nodosalida}   & \fn{nodosalida}   & Salida publicada: tabla, figura o mapa \\
\addlinespace
\muestraarista{flujo}           & \fn{flujo}           & Dependencia normal: el destino lee lo que el origen produce \\
\muestraarista{flujoexcepcion}  & \fn{flujoexcepcion}  & Dependencia de un insumo \textbf{curado}: entra al flujo sin pasar por ningún script \\
\muestraarista{paso}            & \fn{paso}            & Etapa del flujo en el diagrama de arquitectura \\
\muestraarista{control}         & \fn{control}         & Verificación: lee sin modificar \\
\bottomrule
\end{tabularx}

\vspace{6pt}
La distinción entre \fn{flujo} y \fn{flujoexcepcion} no es decorativa: es el
hallazgo estructural del proyecto. Una arista roja discontinua señala un dato
que entra al informe sin haber pasado por código, y por lo tanto sin que exista
forma de regenerarlo.

\subsection{Formato de cifras, fechas y territorios}

\begin{itemize}[leftmargin=*, itemsep=2pt]
  \item \textbf{Números:} convención colombiana. Punto para miles, coma decimal:
    \fn{1.234,5}. Los porcentajes llevan un decimal y espacio antes del signo.
  \item \textbf{Fechas:} \fn{AAAA-MM-DD} en el texto y en los CSV de auditoría.
    En los nombres de archivo de las fuentes, \fn{AAAAMMDD} sin separadores,
    porque así los entrega el equipo.
  \item \textbf{Años de referencia:} un indicador se cita siempre con su año.
    Distintos indicadores tienen cortes distintos y algunos son posteriores al
    informe de 2026.
  \item \textbf{Municipios:} se nombran con su forma de presentación —con tildes
    y preposiciones en minúscula— y se identifican por su código DIVIPOLA.
  \item \textbf{Unidades:} van en el encabezado de la columna o en el subtítulo
    de la figura, nunca solo en el texto que la acompaña.
\end{itemize}

% =============================================================================
\section{El grafo de procesamiento}
\label{sec:dag}
% =============================================================================

\subsection{Arquitectura}

El flujo tiene cuatro etapas y dos almacenes. La dirección es siempre la misma y
no hay ciclos: ninguna etapa escribe sobre lo que otra ya leyó.

\begin{figure}[htbp]\centering
\begin{tikzpicture}[x=1cm, y=1cm]
  % Fila 1: la cadena reproducible, de la fuente a la tabla.
  \node[almacen] (fuentes)   at (0,0)
    {\textbf{Fuentes}\\[2pt]\texttt{\scriptsize 00\_Inputs}\\ \cifraInsumos{} archivos};
  \node[etapa]   (homog)     at (4,0)
    {\textbf{Homogeneización}\\[2pt]\texttt{\scriptsize 01\_homogeneizacion}\\ 11 scripts};
  \node[almacen] (derivados) at (8,0)
    {\textbf{Derivados}\\[2pt]\texttt{\scriptsize 01\_Derived}\\ \cifraDerivados{} tablas Parquet};
  \node[etapa]   (tablas)    at (12,0)
    {\textbf{Tablas}\\[2pt]\texttt{\scriptsize 03\_tablas}\\ 10 secciones};

  % Fila 2: de la tabla a la pieza publicada.
  \node[etapa]  (figuras) at (8,-2.7)
    {\textbf{Figuras y mapas}\\[2pt]\texttt{\scriptsize 04\_figuras}\\ \cifraFiguras{} archivos};
  \node[salida] (salidas) at (12,-2.7)
    {\textbf{Salidas}\\[2pt]\texttt{\scriptsize 03\_Outputs}\\ .xlsx $\cdot$ .png $\cdot$ .pdf};
  \node[almacen] (checks) at (8,-5.0)
    {\textbf{Verificación}\\[2pt]\texttt{\scriptsize 99\_checks}\\ 5 controles};

  \draw[paso] (fuentes)   -- (homog);
  \draw[paso] (homog)     -- (derivados);
  \draw[paso] (derivados) -- (tablas);
  \draw[paso] (tablas.south) -- (salidas.north);
  \draw[paso] (tablas.south west) -- (figuras.north east);
  \draw[paso] (figuras)   -- (salidas);
  \draw[control] (checks.east) -| (salidas.south);

  % Los insumos curados saltan la homogeneización y entran directo a las tablas.
  \draw[flujoexcepcion, line width=0.7pt]
    (fuentes.north) .. controls +(0,1.15) and +(0,1.15) .. (tablas.north);
  \node[font=\scriptsize, text=baya, align=center] at (6,1.95)
    {\cifraCurados{} insumos curados entran sin pasar por ningún script};
\end{tikzpicture}
\caption{Arquitectura del flujo. Las cifras salen del propio repositorio; la
arista roja discontinua es la ruta que no tiene código detrás.}
\end{figure}

\begin{regla}[title={Regla 1 — la figura no recalcula}]
Una figura y un mapa leen la tabla que la sección acaba de escribir. Si
recalcularan desde el derivado podrían discrepar de ella, y el informe tendría
dos cifras distintas para el mismo hecho. Por eso los mapas corren después de
las diez secciones y no dentro de ellas.
\end{regla}

\subsection{Cadena de homogeneización}

Qué script produce qué derivado. El grafo se deriva de leer el código, no de la
documentación: un diagrama dibujado a mano envejece en silencio, este no puede
discrepar del pipeline.

\begin{figure}[htbp]\centering
% Generado por 02_Code/99_checks/generar_tablas_anexo.R — no editar a mano.
\begin{tikzpicture}[x=1cm, y=1cm]
\node[nodoscript, text width=5.2cm] (i1) at (0,-1.440) {crosswalk\_\allowbreak{}territorial.R};
\node[nodoscript, text width=5.2cm] (i2) at (0,-2.160) {deyc.R};
\node[nodoscript, text width=5.2cm] (i3) at (0,-2.880) {ecv\_\allowbreak{}oficial\_\allowbreak{}agregados.R};
\node[nodoscript, text width=5.2cm] (i4) at (0,-3.600) {ecv.R};
\node[nodoscript, text width=5.2cm] (i5) at (0,-4.320) {educacion.R};
\node[nodoscript, text width=5.2cm] (i6) at (0,-5.040) {infraestructura\_\allowbreak{}internet.R};
\node[nodoscript, text width=5.2cm] (i7) at (0,-5.760) {percepcion\_\allowbreak{}seguridad.R};
\node[nodoscript, text width=5.2cm] (i8) at (0,-6.480) {poblacion\_\allowbreak{}municipal.R};
\node[nodoscript, text width=5.2cm] (i9) at (0,-7.200) {salud\_\allowbreak{}seguridad.R};
\node[nodoscript, text width=5.2cm] (i10) at (0,-7.920) {seguridad\_\allowbreak{}policia.R};
\node[nodoscript, text width=5.2cm] (i11) at (0,-8.640) {uso\_\allowbreak{}suelo.R};
\node[nododerivado, text width=6.8cm] (d1) at (8.8,-0.000) {codigos\_\allowbreak{}provincias};
\node[nododerivado, text width=6.8cm] (d2) at (8.8,-0.720) {subreg\_\allowbreak{}completo};
\node[nododerivado, text width=6.8cm] (d3) at (8.8,-1.440) {estructura\_\allowbreak{}productiva};
\node[nododerivado, text width=6.8cm] (d4) at (8.8,-2.160) {imca\_\allowbreak{}2022};
\node[nododerivado, text width=6.8cm] (d5) at (8.8,-2.880) {ECV\_\allowbreak{}oficial\_\allowbreak{}agregados};
\node[nododerivado, text width=6.8cm] (d6) at (8.8,-3.600) {ECV\_\allowbreak{}nbi\_\allowbreak{}pobreza};
\node[nododerivado, text width=6.8cm] (d7) at (8.8,-4.320) {ECV\_\allowbreak{}raw};
\node[nododerivado, text width=6.8cm] (d8) at (8.8,-5.040) {20250506 EDUCACION};
\node[nododerivado, text width=6.8cm] (d9) at (8.8,-5.760) {infraestructura\_\allowbreak{}internet\_\allowbreak{}2025};
\node[nododerivado, text width=6.8cm] (d10) at (8.8,-6.480) {percepcion\_\allowbreak{}seguridad};
\node[nododerivado, text width=6.8cm] (d11) at (8.8,-7.200) {poblacion\_\allowbreak{}municipal\_\allowbreak{}total\_\allowbreak{}2025};
\node[nododerivado, text width=6.8cm] (d12) at (8.8,-7.920) {poblacion\_\allowbreak{}total\_\allowbreak{}2025};
\node[nododerivado, text width=6.8cm] (d13) at (8.8,-8.640) {20260504\_\allowbreak{}SEGURIDAD\_\allowbreak{}SALUD\_\allowbreak{}DEFICT\_\allowbreak{}VIVIENDA};
\node[nododerivado, text width=6.8cm] (d14) at (8.8,-9.360) {seguridad\_\allowbreak{}policia};
\node[nododerivado, text width=6.8cm] (d15) at (8.8,-10.080) {uso\_\allowbreak{}del\_\allowbreak{}suelo\_\allowbreak{}pota};
\draw[flujo] (i1.east) -- (d1.west);
\draw[flujo] (i1.east) -- (d2.west);
\draw[flujo] (i2.east) -- (d3.west);
\draw[flujo] (i2.east) -- (d4.west);
\draw[flujo] (i3.east) -- (d5.west);
\draw[flujo] (i4.east) -- (d7.west);
\draw[flujo] (i4.east) -- (d6.west);
\draw[flujo] (i5.east) -- (d8.west);
\draw[flujo] (i6.east) -- (d9.west);
\draw[flujo] (i7.east) -- (d10.west);
\draw[flujo] (i8.east) -- (d12.west);
\draw[flujo] (i8.east) -- (d11.west);
\draw[flujo] (i9.east) -- (d13.west);
\draw[flujo] (i10.east) -- (d14.west);
\draw[flujo] (i11.east) -- (d15.west);
\end{tikzpicture}

\caption{De script de homogeneización a derivado. Los
\cifraDerivados{} derivados tienen productor identificado; cuatro scripts
producen más de uno. Este grafo cubre solo la ruta reproducible: los insumos
curados no aparecen porque ningún script los escribe.}
\end{figure}

\subsection{Consumo por sección}

Qué lee cada sección temática. Es el grafo que importa para responder «¿de dónde
salió esta cifra?», y el que muestra que las secciones se alimentan de tres
sitios distintos: derivados del pipeline, tableros curados y fuentes crudas
leídas directamente.

\begin{figure}[p]\centering
% Generado por 02_Code/99_checks/generar_tablas_anexo.R — no editar a mano.
\begin{tikzpicture}[x=1cm, y=1cm]
\node[nodorecurso, text width=7.8cm] (i1) at (0,-0.000) {20260504\_\allowbreak{}SEGURIDAD\_\allowbreak{}SALUD\_\allowbreak{}DEFICT\_\allowbreak{}VIVIENDA};
\node[nodorecurso, text width=7.8cm] (i2) at (0,-0.620) {AREA\_\allowbreak{}ORIGINAL.xlsx};
\node[nodorecurso, text width=7.8cm] (i3) at (0,-1.240) {INDICADORES ECV 2023 MUNICIPIOS.xlsx};
\node[nodorecurso, text width=7.8cm] (i4) at (0,-1.860) {poblacion\_\allowbreak{}municipal\_\allowbreak{}total\_\allowbreak{}2025\_\allowbreak{}dicc.dta};
\node[nodorecurso, text width=7.8cm] (i5) at (0,-2.480) {poblacion\_\allowbreak{}total\_\allowbreak{}2025};
\node[nodorecurso, text width=7.8cm] (i6) at (0,-3.100) {20250506 EDUCACION};
\node[nodorecurso, text width=7.8cm] (i7) at (0,-3.720) {CATASTRO.xlsx};
\node[nodorecurso, text width=7.8cm] (i8) at (0,-4.340) {deficit cualitativo por municipio 2023.xlsx};
\node[nodorecurso, text width=7.8cm] (i9) at (0,-4.960) {deficit cuantitativo por municipio 2023.xlsx};
\node[nodorecurso, text width=7.8cm] (i10) at (0,-5.580) {ESTADO\_\allowbreak{}CATASTRO.xlsx};
\node[nodorecurso, text width=7.8cm] (i11) at (0,-6.200) {imca\_\allowbreak{}2022};
\node[nodorecurso, text width=7.8cm] (i12) at (0,-6.820) {indice\_\allowbreak{}gobierno\_\allowbreak{}digital\_\allowbreak{}2024.dta};
\node[nodorecurso, text width=7.8cm] (i13) at (0,-7.440) {infraestructura\_\allowbreak{}internet\_\allowbreak{}2025};
\node[nodorecurso, text width=7.8cm] (i14) at (0,-8.060) {infraestructura\_\allowbreak{}municipios\_\allowbreak{}dicc.dta};
\node[nodorecurso, text width=7.8cm] (i15) at (0,-8.680) {uso\_\allowbreak{}del\_\allowbreak{}suelo\_\allowbreak{}pota};
\node[nodorecurso, text width=7.8cm] (i16) at (0,-9.300) {POBLACION MUNICIPAL.xlsx};
\node[nodorecurso, text width=7.8cm] (i17) at (0,-9.920) {617\_\allowbreak{}PROVINCIAS.xlsx};
\node[nodorecurso, text width=7.8cm] (i18) at (0,-10.540) {ECV\_\allowbreak{}nbi\_\allowbreak{}pobreza};
\node[nodorecurso, text width=7.8cm] (i19) at (0,-11.160) {ECV\_\allowbreak{}oficial\_\allowbreak{}agregados};
\node[nodorecurso, text width=7.8cm] (i20) at (0,-11.780) {ICM\_\allowbreak{}PROVINCIAS.xlsx};
\node[nodorecurso, text width=7.8cm] (i21) at (0,-12.400) {IDF\_\allowbreak{}PROVINCIAS.xlsx};
\node[nodorecurso, text width=7.8cm] (i22) at (0,-13.020) {MDM\_\allowbreak{}PROVINCIAS.xlsx};
\node[nodorecurso, text width=7.8cm] (i23) at (0,-13.640) {Actualización/Fuentes/TITULOSMINAS\_\allowbreak{}ORIGINAL.xlsx};
\node[nodorecurso, text width=7.8cm] (i24) at (0,-14.260) {DATALAKE DEyC.xlsx};
\node[nodorecurso, text width=7.8cm] (i25) at (0,-14.880) {Energia\_\allowbreak{}NorteAntioquia.xlsx};
\node[nodorecurso, text width=7.8cm] (i26) at (0,-15.500) {EXTRANJEROS\_\allowbreak{}PROVINCIAS.xlsx};
\node[nodorecurso, text width=7.8cm] (i27) at (0,-16.120) {instrumentos\_\allowbreak{}planificacion\_\allowbreak{}turistica\_\allowbreak{}antioquia.xlsx};
\node[nodorecurso, text width=7.8cm] (i28) at (0,-16.740) {VALOR\_\allowbreak{}AGREGADO\_\allowbreak{}20152024\_\allowbreak{}MUNICIPIOS.xlsx};
\node[nodoseccion, text width=4.2cm] (d1) at (11.4,-7.130) {01 generalidades};
\node[nodoseccion, text width=4.2cm] (d2) at (11.4,-7.750) {02 demografia};
\node[nodoseccion, text width=4.2cm] (d3) at (11.4,-8.370) {03 ordenamiento};
\node[nodoseccion, text width=4.2cm] (d4) at (11.4,-8.990) {04 gobernabilidad};
\node[nodoseccion, text width=4.2cm] (d5) at (11.4,-9.610) {05 economia};
\draw[flujo] (i1.east) -- (d2.west);
\draw[flujoexcepcion] (i2.east) -- (d1.west);
\draw[flujoexcepcion] (i2.east) -- (d2.west);
\draw[flujoexcepcion] (i2.east) -- (d3.west);
\draw[flujoexcepcion] (i3.east) -- (d2.west);
\draw[flujoexcepcion] (i4.east) -- (d2.west);
\draw[flujo] (i5.east) -- (d2.west);
\draw[flujo] (i5.east) -- (d3.west);
\draw[flujo] (i6.east) -- (d3.west);
\draw[flujoexcepcion] (i7.east) -- (d3.west);
\draw[flujoexcepcion] (i8.east) -- (d3.west);
\draw[flujoexcepcion] (i9.east) -- (d3.west);
\draw[flujoexcepcion] (i10.east) -- (d3.west);
\draw[flujo] (i11.east) -- (d3.west);
\draw[flujoexcepcion] (i12.east) -- (d3.west);
\draw[flujo] (i13.east) -- (d3.west);
\draw[flujoexcepcion] (i14.east) -- (d3.west);
\draw[flujo] (i15.east) -- (d3.west);
\draw[flujoexcepcion] (i16.east) -- (d2.west);
\draw[flujoexcepcion] (i16.east) -- (d3.west);
\draw[flujoexcepcion] (i16.east) -- (d4.west);
\draw[flujoexcepcion] (i16.east) -- (d5.west);
\draw[flujoexcepcion] (i17.east) -- (d4.west);
\draw[flujo] (i18.east) -- (d3.west);
\draw[flujo] (i18.east) -- (d5.west);
\draw[flujo] (i19.east) -- (d3.west);
\draw[flujo] (i19.east) -- (d5.west);
\draw[flujoexcepcion] (i20.east) -- (d4.west);
\draw[flujoexcepcion] (i21.east) -- (d4.west);
\draw[flujoexcepcion] (i22.east) -- (d4.west);
\draw[flujoexcepcion] (i23.east) -- (d5.west);
\draw[flujoexcepcion] (i24.east) -- (d5.west);
\draw[flujoexcepcion] (i25.east) -- (d5.west);
\draw[flujoexcepcion] (i26.east) -- (d5.west);
\draw[flujoexcepcion] (i27.east) -- (d5.west);
\draw[flujoexcepcion] (i28.east) -- (d5.west);
\end{tikzpicture}

\caption{Recursos que consumen las secciones 01 a 05. En rojo discontinuo, los
insumos curados: entran al informe sin pasar por ningún script.}
\end{figure}

\begin{figure}[p]\centering
% Generado por 02_Code/99_checks/generar_tablas_anexo.R — no editar a mano.
\begin{tikzpicture}[x=1cm, y=1cm]
\node[nodorecurso, text width=7.8cm] (i1) at (0,-0.000) {CULTIVOS\_\allowbreak{}PROVINCIAS.xlsx};
\node[nodorecurso, text width=7.8cm] (i2) at (0,-0.620) {PECUARIO\_\allowbreak{}PROVINCIAS.xlsx};
\node[nodorecurso, text width=7.8cm] (i3) at (0,-1.240) {AREAPROTEGIDA\_\allowbreak{}PROVINCIAS.xlsx};
\node[nodorecurso, text width=7.8cm] (i4) at (0,-1.860) {EMERGENCIAS\_\allowbreak{}PROVINCIAS.xlsx};
\node[nodorecurso, text width=7.8cm] (i5) at (0,-2.480) {IMRC\_\allowbreak{}PROVINCIAS.xlsx};
\node[nodorecurso, text width=7.8cm] (i6) at (0,-3.100) {IRCA\_\allowbreak{}PROVINCIAS.xlsx};
\node[nodorecurso, text width=7.8cm] (i7) at (0,-3.720) {Pérdida de cobertura árborea.xlsx};
\node[nodorecurso, text width=7.8cm] (i8) at (0,-4.340) {AREA\_\allowbreak{}ORIGINAL.xlsx};
\node[nodorecurso, text width=7.8cm] (i9) at (0,-4.960) {EDUCACION\_\allowbreak{}PROVINCIAS.xlsx};
\node[nodorecurso, text width=7.8cm] (i10) at (0,-5.580) {20260504\_\allowbreak{}SEGURIDAD\_\allowbreak{}SALUD\_\allowbreak{}DEFICT\_\allowbreak{}VIVIENDA};
\node[nodorecurso, text width=7.8cm] (i11) at (0,-6.200) {ASEGURAMIENTO\_\allowbreak{}MUNICIPIOS.xlsx};
\node[nodorecurso, text width=7.8cm] (i12) at (0,-6.820) {MORTALIDAD\_\allowbreak{}INFANTIL.xlsx};
\node[nodorecurso, text width=7.8cm] (i13) at (0,-7.440) {NATALIDAD.xlsx};
\node[nodorecurso, text width=7.8cm] (i14) at (0,-8.060) {POBLACION MUNICIPAL.xlsx};
\node[nodorecurso, text width=7.8cm] (i15) at (0,-8.680) {suicidios\_\allowbreak{}e\_\allowbreak{}intentos\_\allowbreak{}medias\_\allowbreak{}dicc.dta};
\node[nodorecurso, text width=7.8cm] (i16) at (0,-9.300) {CULTILICITOS\_\allowbreak{}PROVINCIAS.xlsx};
\node[nodorecurso, text width=7.8cm] (i17) at (0,-9.920) {EVOA\_\allowbreak{}PROVINCIAS.xlsx};
\node[nodorecurso, text width=7.8cm] (i18) at (0,-10.540) {HECHOSVICTIM\_\allowbreak{}PROVINCIAS.xlsx};
\node[nodorecurso, text width=7.8cm] (i19) at (0,-11.160) {IRV-2025.xlsx};
\node[nodorecurso, text width=7.8cm] (i20) at (0,-11.780) {percepcion\_\allowbreak{}seguridad};
\node[nodorecurso, text width=7.8cm] (i21) at (0,-12.400) {poblacion\_\allowbreak{}total\_\allowbreak{}2025};
\node[nodorecurso, text width=7.8cm] (i22) at (0,-13.020) {seguridad\_\allowbreak{}policia};
\node[nodorecurso, text width=7.8cm] (i23) at (0,-13.640) {SRDAFT\_\allowbreak{}PROVINCIAS.xlsx};
\node[nodorecurso, text width=7.8cm] (i24) at (0,-14.260) {UBPD\_\allowbreak{}desaparecidos\_\allowbreak{}Antioquia.xlsx};
\node[nodoseccion, text width=4.2cm] (d1) at (11.4,-5.890) {06 desarrollo rural};
\node[nodoseccion, text width=4.2cm] (d2) at (11.4,-6.510) {07 ambiental};
\node[nodoseccion, text width=4.2cm] (d3) at (11.4,-7.130) {08 educacion};
\node[nodoseccion, text width=4.2cm] (d4) at (11.4,-7.750) {09 salud};
\node[nodoseccion, text width=4.2cm] (d5) at (11.4,-8.370) {10 seguridad};
\draw[flujoexcepcion] (i1.east) -- (d1.west);
\draw[flujoexcepcion] (i2.east) -- (d1.west);
\draw[flujoexcepcion] (i3.east) -- (d2.west);
\draw[flujoexcepcion] (i4.east) -- (d2.west);
\draw[flujoexcepcion] (i5.east) -- (d2.west);
\draw[flujoexcepcion] (i6.east) -- (d2.west);
\draw[flujoexcepcion] (i7.east) -- (d2.west);
\draw[flujoexcepcion] (i8.east) -- (d1.west);
\draw[flujoexcepcion] (i8.east) -- (d5.west);
\draw[flujoexcepcion] (i9.east) -- (d3.west);
\draw[flujo] (i10.east) -- (d4.west);
\draw[flujoexcepcion] (i11.east) -- (d4.west);
\draw[flujoexcepcion] (i12.east) -- (d4.west);
\draw[flujoexcepcion] (i13.east) -- (d4.west);
\draw[flujoexcepcion] (i14.east) -- (d4.west);
\draw[flujoexcepcion] (i15.east) -- (d4.west);
\draw[flujoexcepcion] (i16.east) -- (d5.west);
\draw[flujoexcepcion] (i17.east) -- (d5.west);
\draw[flujoexcepcion] (i18.east) -- (d5.west);
\draw[flujoexcepcion] (i19.east) -- (d5.west);
\draw[flujo] (i20.east) -- (d5.west);
\draw[flujo] (i21.east) -- (d5.west);
\draw[flujo] (i22.east) -- (d5.west);
\draw[flujoexcepcion] (i23.east) -- (d5.west);
\draw[flujoexcepcion] (i24.east) -- (d5.west);
\end{tikzpicture}

\caption{Recursos que consumen las secciones 06 a 10.}
\end{figure}

\clearpage

% =============================================================================
\section{Fuentes de datos}
% =============================================================================

\subsection{Catálogo}

El código referencia \cifraInsumos{} archivos de entrada, de los cuales
\cifraInsumosPresentes{} están en el repositorio y \cifraInsumosAusentes{} deben
pedirse al equipo. La tabla~\ref{tab:fuentes} los lista con su peso, su fecha de
última modificación y cuántos scripts los leen.

\begin{table}[htbp]\centering
\caption{Catálogo de fuentes que el flujo lee}
\label{tab:fuentes}
\scriptsize
% Generado por 02_Code/99_checks/generar_tablas_anexo.R — no editar a mano.
\begin{tabularx}{\linewidth}{>{\raggedright\arraybackslash}X l r l r}
\toprule
\textbf{Archivo} & \textbf{Tipo} & \textbf{Peso} & \textbf{Última modificación} & \textbf{Scripts} \\
\midrule
MGN\_\allowbreak{}MPIO\_\allowbreak{}POLITICO.shp & Actualización & 9 Mb & 2026-08-08 06:47 & 1 \\
TITULOSMINAS\_\allowbreak{}ORIGINAL.xlsx & Actualización & 947 Kb & 2026-08-08 06:52 & 1 \\
AREA\_\allowbreak{}ORIGINAL.xlsx & Crudo & 12 Kb & 2026-08-08 06:52 & 5 \\
ASEGURAMIENTO\_\allowbreak{}MUNICIPIOS.xlsx & Crudo & 79 Kb & 2026-08-08 06:52 & 1 \\
CATASTRO.xlsx & Crudo & 40 Kb & 2026-08-08 06:52 & 1 \\
códigos\_\allowbreak{}municipios\_\allowbreak{}clean.dta & Crudo & 14 Kb & 2026-08-08 06:52 & 1 \\
DATALAKE DEyC.xlsx & Crudo & 2 Mb & 2026-08-08 06:47 & 2 \\
DATALAKE EDUCACION.xlsx & Crudo & 7 Mb & 2026-08-08 06:52 & 1 \\
DATOS SALUD.xlsx & Crudo & 15 Mb & 2026-08-08 06:52 & 1 \\
deficit cualitativo por municipio 2023.xlsx & Crudo & 71 Kb & 2026-08-08 06:47 & 1 \\
deficit cuantitativo por municipio 2023.xlsx & Crudo & 48 Kb & 2026-08-08 06:47 & 1 \\
EMPAQUETAMIENTO\_\allowbreak{}FIJO\_\allowbreak{}3.csv & Crudo & — & no versionado & 1 \\
Energia\_\allowbreak{}NorteAntioquia.xlsx & Crudo & 15 Kb & 2026-08-08 06:52 & 1 \\
INDICADORES ECV 2023 MUNICIPIOS.xlsx & Crudo & 28 Mb & 2026-08-08 06:47 & 2 \\
Indicadores\_\allowbreak{}ECV2023.xlsx & Crudo & — & no versionado & 1 \\
instrumentos\_\allowbreak{}planificacion\_\allowbreak{}turistica\_\allowbreak{}antioquia.xlsx & Crudo & 26 Kb & 2026-08-08 06:52 & 1 \\
IRV-2025.xlsx & Crudo & 439 Kb & 2026-08-08 06:52 & 1 \\
MORTALIDAD\_\allowbreak{}INFANTIL.xlsx & Crudo & 375 Kb & 2026-08-08 06:52 & 1 \\
MUNICIPIOS\_\allowbreak{}SUBREG\_\allowbreak{}PROV.xlsx & Crudo & 13 Kb & 2026-08-08 06:52 & 1 \\
NATALIDAD.xlsx & Crudo & 514 Kb & 2026-08-08 06:52 & 2 \\
Pérdida de cobertura árborea.xlsx & Crudo & 198 Kb & 2026-08-08 06:47 & 1 \\
POBLACION MUNICIPAL.xlsx & Crudo & 36 Mb & 2026-08-08 06:47 & 5 \\
UBPD\_\allowbreak{}desaparecidos\_\allowbreak{}Antioquia.xlsx & Crudo & 26 Kb & 2026-08-08 06:52 & 1 \\
617\_\allowbreak{}PROVINCIAS.xlsx & Curado & 91 Kb & 2026-08-08 06:52 & 1 \\
AREAPROTEGIDA\_\allowbreak{}PROVINCIAS.xlsx & Curado & 24 Kb & 2026-08-08 06:52 & 1 \\
CULTILICITOS\_\allowbreak{}PROVINCIAS.xlsx & Curado & 11 Kb & 2026-08-08 06:52 & 1 \\
CULTIVOS\_\allowbreak{}PROVINCIAS.xlsx & Curado & 839 Kb & 2026-08-08 06:52 & 1 \\
EDUCACION\_\allowbreak{}PROVINCIAS.xlsx & Curado & 108 Kb & 2026-08-08 06:52 & 1 \\
EMERGENCIAS\_\allowbreak{}PROVINCIAS.xlsx & Curado & 127 Kb & 2026-08-08 06:52 & 1 \\
ESTADO\_\allowbreak{}CATASTRO.xlsx & Curado & 14 Kb & 2026-08-08 06:52 & 1 \\
EVOA\_\allowbreak{}PROVINCIAS.xlsx & Curado & 20 Kb & 2026-08-08 06:52 & 1 \\
EXTRANJEROS\_\allowbreak{}PROVINCIAS.xlsx & Curado & 57 Kb & 2026-08-08 06:52 & 1 \\
HECHOSVICTIM\_\allowbreak{}PROVINCIAS.xlsx & Curado & 82 Kb & 2026-08-08 06:52 & 1 \\
ICM\_\allowbreak{}PROVINCIAS.xlsx & Curado & 325 Kb & 2026-08-08 06:52 & 1 \\
IDF\_\allowbreak{}PROVINCIAS.xlsx & Curado & 19 Kb & 2026-08-08 06:52 & 1 \\
IMRC\_\allowbreak{}PROVINCIAS.xlsx & Curado & 19 Kb & 2026-08-08 06:52 & 1 \\
indice\_\allowbreak{}gobierno\_\allowbreak{}digital\_\allowbreak{}2024.dta & Curado & 25 Kb & 2026-08-08 06:52 & 1 \\
infraestructura\_\allowbreak{}municipios\_\allowbreak{}dicc.dta & Curado & 27 Kb & 2026-08-08 06:47 & 1 \\
IRCA\_\allowbreak{}PROVINCIAS.xlsx & Curado & 130 Kb & 2026-08-08 06:52 & 1 \\
MDM\_\allowbreak{}PROVINCIAS.xlsx & Curado & 107 Kb & 2026-08-08 06:52 & 1 \\
PECUARIO\_\allowbreak{}PROVINCIAS.xlsx & Curado & 5 Mb & 2026-08-08 06:52 & 1 \\
poblacion\_\allowbreak{}municipal\_\allowbreak{}total\_\allowbreak{}2025\_\allowbreak{}dicc.dta & Curado & 151 Kb & 2026-08-08 06:52 & 1 \\
SRDAFT\_\allowbreak{}PROVINCIAS.xlsx & Curado & 16 Kb & 2026-08-08 06:52 & 1 \\
suicidios\_\allowbreak{}e\_\allowbreak{}intentos\_\allowbreak{}medias\_\allowbreak{}dicc.dta & Curado & 38 Kb & 2026-08-08 06:47 & 1 \\
VALOR\_\allowbreak{}AGREGADO\_\allowbreak{}20152024\_\allowbreak{}MUNICIPIOS.xlsx & Curado & 512 Kb & 2026-08-08 06:52 & 1 \\
\bottomrule
\addlinespace[2pt]
\multicolumn{5}{@{}p{\linewidth}@{}}{\footnotesize «Crudo»: fuente original, solo lectura. «Curado»: tablero que el equipo arma a mano en Excel y ningún script regenera. «Scripts»: cuántos scripts del pipeline leen el archivo. El listado con la huella MD5 de cada archivo está en 04\_Docs/auditoria/B1\_inventario\_insumos.csv.} \\
\end{tabularx}

\end{table}

\subsection{Las tres clases de insumo}

\begin{description}[leftmargin=*, style=nextline, itemsep=3pt]
  \item[Crudo] Fuente original de una entidad (DANE, Gobernación de Antioquia,
    Policía Nacional, MinSalud, MinTIC, DNP, UPRA, ANM, XM, UNGRD, UBPD, Global
    Forest Watch). Se reemplaza conservando el nombre del archivo.
  \item[Actualización] Fuentes nuevas o revisadas que llegaron después del corte
    original. Viven en \ruta{00\_Inputs/Actualización/Fuentes/} y algunas ya son
    parte del flujo vigente: la carpeta no es un borrador.
  \item[Curado] Tableros que el equipo arma a mano en Excel y que ningún script
    regenera. Son \cifraCurados{} archivos y constituyen el mayor riesgo de
    reproducibilidad del proyecto (sección~\ref{sec:limitaciones}). Viven
    separados en \ruta{00\_Inputs/curados/} precisamente para que no se
    confundan con lo que sí se regenera.
\end{description}

\begin{table}[htbp]\centering
\caption{Fuentes que no viajan con el repositorio}
\small
% Generado por 02_Code/99_checks/generar_tablas_anexo.R — no editar a mano.
\begin{tabularx}{\linewidth}{>{\raggedright\arraybackslash}X X}
\toprule
\textbf{Archivo} & \textbf{Script que lo necesita} \\
\midrule
Indicadores\_\allowbreak{}ECV2023.xlsx & 02\_\allowbreak{}Code/01\_\allowbreak{}homogeneizacion/ecv\_\allowbreak{}oficial\_\allowbreak{}agregados.R \\
EMPAQUETAMIENTO\_\allowbreak{}FIJO\_\allowbreak{}3.csv & 02\_\allowbreak{}Code/01\_\allowbreak{}homogeneizacion/infraestructura\_\allowbreak{}internet.R \\
\bottomrule
\addlinespace[2pt]
\multicolumn{2}{@{}p{\linewidth}@{}}{\footnotesize Estas fuentes no caben en el repositorio o contienen microdatos que no se versionan. El pipeline sigue corriendo sin ellas: la etapa que las usa se detiene con un mensaje que dice cuál falta, y el resto continúa.} \\
\end{tabularx}

\end{table}

\subsection{Derivados}

La homogeneización produce \cifraDerivados{} derivados, todos en formato
Parquet: columnar, comprimido, tipado y legible desde R, Python y Stata~18+.
En el código no se abren por ruta sino con \fn{leer\_derivado("nombre")}, de modo
que cambiar el formato mañana no obligue a tocar treinta scripts.

\begin{table}[htbp]\centering
\caption{Derivados del flujo y el script que los produce}
\small
% Generado por 02_Code/99_checks/generar_tablas_anexo.R — no editar a mano.
\begin{tabularx}{\linewidth}{>{\raggedright\arraybackslash}X l r}
\toprule
\textbf{Derivado} & \textbf{Script que lo produce} & \textbf{Consumidores} \\
\midrule
20250506 EDUCACION & educacion.R & 1 \\
20260504\_\allowbreak{}SEGURIDAD\_\allowbreak{}SALUD\_\allowbreak{}DEFICT\_\allowbreak{}VIVIENDA & salud\_\allowbreak{}seguridad.R & 2 \\
codigos\_\allowbreak{}provincias & crosswalk\_\allowbreak{}territorial.R & 1 \\
ECV\_\allowbreak{}nbi\_\allowbreak{}pobreza & ecv.R & 2 \\
ECV\_\allowbreak{}oficial\_\allowbreak{}agregados & ecv\_\allowbreak{}oficial\_\allowbreak{}agregados.R & 2 \\
ECV\_\allowbreak{}raw & ecv.R & 0 \\
estructura\_\allowbreak{}productiva & deyc.R & 0 \\
imca\_\allowbreak{}2022 & deyc.R & 1 \\
infraestructura\_\allowbreak{}internet\_\allowbreak{}2025 & infraestructura\_\allowbreak{}internet.R & 1 \\
percepcion\_\allowbreak{}seguridad & percepcion\_\allowbreak{}seguridad.R & 1 \\
poblacion\_\allowbreak{}municipal\_\allowbreak{}total\_\allowbreak{}2025 & poblacion\_\allowbreak{}municipal.R & 0 \\
poblacion\_\allowbreak{}total\_\allowbreak{}2025 & poblacion\_\allowbreak{}municipal.R & 4 \\
seguridad\_\allowbreak{}policia & seguridad\_\allowbreak{}policia.R & 1 \\
subreg\_\allowbreak{}completo & crosswalk\_\allowbreak{}territorial.R & 1 \\
uso\_\allowbreak{}del\_\allowbreak{}suelo\_\allowbreak{}pota & uso\_\allowbreak{}suelo.R & 1 \\
\bottomrule
\addlinespace[2pt]
\multicolumn{3}{@{}p{\linewidth}@{}}{\footnotesize Todos los derivados se guardan en Parquet en 01\_Data/01\_Derived y se leen con leer\_derivado().} \\
\end{tabularx}

\end{table}

% =============================================================================
\section{Estructura de datos}
\label{sec:estructura}
% =============================================================================

Esta sección declara el modelo: qué representa una fila, qué la identifica, qué
tipo tiene cada columna, cómo se escribe un dato ausente y en qué unidades está
todo. Es lo que hay que leer antes de cruzar estos datos con otros.

\subsection{Modelo y unidad de observación}

El modelo es deliberadamente plano. No hay una base relacional: hay
\emph{tablas municipales} que se cruzan por código DANE contra un único
crosswalk territorial. Tres niveles:

\begin{description}[leftmargin=*, style=nextline, itemsep=3pt]
  \item[Nivel 0 — territorio] \fn{codigos\_provincias} y \fn{subreg\_completo}.
    Definen qué municipio existe y a qué provincia y subregión pertenece. Son la
    tabla maestra: cualquier discrepancia entre fuentes se resuelve a su favor.
  \item[Nivel 1 — indicadores municipales] Los \cifraDerivados{} derivados. Una
    fila por municipio, o por municipio y año cuando la serie es temporal.
  \item[Nivel 2 — salidas por provincia] Las hojas de
    \ruta{03\_Outputs}. Una fila por municipio de la provincia, más las filas
    agregadas de provincia, subregión y departamento, que se distinguen por
    tener el código DANE vacío.
\end{description}

\begin{regla}[title={Regla 8 — la fila agregada se reconoce por el código vacío}]
En toda hoja publicada, una fila sin código DANE es un agregado, no un
municipio. Es la convención de la que dependen los formatos, las figuras y la
auditoría; romperla haría que un total se contara como un territorio más.
\end{regla}

\subsection{Llaves e integridad referencial}

\begin{table}[htbp]\centering
\caption{Llaves del modelo, su dominio y su fuente de verdad}
\small
% Generado por 02_Code/99_checks/generar_tablas_anexo.R — no editar a mano.
\begin{tabularx}{\linewidth}{l l X >{\raggedright\arraybackslash}p{2.9cm} >{\raggedright\arraybackslash}p{3.5cm}}
\toprule
\textbf{Llave} & \textbf{Tipo} & \textbf{Dominio} & \textbf{Cardinalidad} & \textbf{Fuente de verdad} \\
\midrule
ind\_\allowbreak{}mpio & entero & 5001--5895 (DIVIPOLA de Antioquia) & 88 en las provincias; 125 en el departamento & códigos\_\allowbreak{}municipios\_\allowbreak{}clean.dta \\
id\_\allowbreak{}provincia & entero & 1--11 & 11 & tabla PROVINCIAS de 00\_\allowbreak{}config.R \\
nvl\_\allowbreak{}label & texto & MAYÚSCULAS sin tildes & 88 & códigos\_\allowbreak{}municipios\_\allowbreak{}clean.dta \\
subregion & texto & nombre de subregión del municipio & 6 & MUNICIPIOS\_\allowbreak{}SUBREG\_\allowbreak{}PROV.xlsx \\
subregion\_\allowbreak{}full & texto & subregión DANE de los 125 municipios & 9 & MUNICIPIOS\_\allowbreak{}SUBREG\_\allowbreak{}PROV.xlsx \\
\bottomrule
\addlinespace[2pt]
\multicolumn{5}{@{}p{\linewidth}@{}}{\footnotesize Todo cruce entre fuentes se hace por ind\_mpio. Los nombres solo se usan cuando la fuente no trae código, y siempre contra el listado oficial.} \\
\end{tabularx}

\end{table}

Las restricciones que el flujo verifica en cada corrida, y que abortan la
ejecución si dejan de cumplirse:

\begin{itemize}[leftmargin=*, itemsep=2pt]
  \item \fn{ind\_mpio} es único en el crosswalk y no admite vacíos.
  \item El crosswalk tiene exactamente \cifraMunicipios{} municipios repartidos
    en las \cifraProvincias{} provincias, y ninguna provincia queda sin
    municipios.
  \item Todo nombre de esquema asociativo del listado maestro empareja con una
    de las \cifraProvincias{} provincias declaradas.
  \item Ninguna hoja publicada contiene municipios ajenos a su provincia
    (verificado en las \cifraProvincias{} provincias, sección~\ref{sec:trazabilidad}).
\end{itemize}

\subsection{Esquema de los derivados}

Las \cifraTablasDerivadas{} tablas derivadas suman \cifraColumnasDerivadas{}
columnas. La llave de cada una no se declara a mano: se busca en cada corrida
probando combinaciones de columnas hasta dar con la más pequeña sin
repeticiones.

\begin{table}[htbp]\centering
\caption{Estructura de las tablas derivadas}
\scriptsize
% Generado por 02_Code/99_checks/generar_tablas_anexo.R — no editar a mano.
\begin{tabularx}{\linewidth}{>{\raggedright\arraybackslash}p{4.9cm} r r >{\raggedright\arraybackslash}p{3.0cm} X}
\toprule
\textbf{Tabla} & \textbf{Filas} & \textbf{Cols.} & \textbf{Llave detectada} & \textbf{Granularidad} \\
\midrule
20250506 EDUCACION & 125 & 18 & ind\_\allowbreak{}mpio & una fila por municipio \\
20260504\_\allowbreak{}SEGURIDAD\_\allowbreak{}SALUD\_\allowbreak{}DEFICT\_\allowbreak{}VIVIENDA & 125 & 36 & ind\_\allowbreak{}mpio & una fila por municipio \\
codigos\_\allowbreak{}provincias & 88 & 6 & ind\_\allowbreak{}mpio & una fila por municipio \\
ECV\_\allowbreak{}nbi\_\allowbreak{}pobreza & 125 & 131 & ind\_\allowbreak{}mpio & una fila por municipio \\
ECV\_\allowbreak{}oficial\_\allowbreak{}agregados & 16 & 51 & key & una fila por territorio \\
ECV\_\allowbreak{}raw & 99250 & 26 & Indicador + Municipio & una fila por indicador × municipio \\
estructura\_\allowbreak{}productiva & 135 & 8 & ind\_\allowbreak{}mpio & una fila por municipio \\
imca\_\allowbreak{}2022 & 134 & 14 & ind\_\allowbreak{}mpio & una fila por municipio \\
infraestructura\_\allowbreak{}internet\_\allowbreak{}2025 & 124190 & 23 & sin llave única & no determinada \\
percepcion\_\allowbreak{}seguridad & 23216 & 8 & sin llave única & no determinada \\
poblacion\_\allowbreak{}municipal\_\allowbreak{}total\_\allowbreak{}2025 & 1123 & 95 & ind\_\allowbreak{}mpio & una fila por municipio \\
poblacion\_\allowbreak{}total\_\allowbreak{}2025 & 375 & 9 & ind\_\allowbreak{}mpio + area\_\allowbreak{}geo & una fila por municipio × área geográfica \\
seguridad\_\allowbreak{}policia & 1000 & 20 & ind\_\allowbreak{}mpio + delito + anio & una fila por municipio × delito × año \\
subreg\_\allowbreak{}completo & 125 & 2 & ind\_\allowbreak{}mpio & una fila por municipio \\
uso\_\allowbreak{}del\_\allowbreak{}suelo\_\allowbreak{}pota & 88 & 17 & codmpio & una fila por municipio \\
\bottomrule
\addlinespace[2pt]
\multicolumn{5}{@{}p{\linewidth}@{}}{\footnotesize La llave no se declara: se busca en cada corrida probando combinaciones de columnas hasta encontrar la más pequeña sin repeticiones. Así deja de ser una promesa y pasa a ser una propiedad verificada.} \\
\end{tabularx}

\end{table}

El esquema columna a columna —tipo, número de valores distintos, porcentaje de
vacíos y un valor de ejemplo— está en el apéndice~\ref{ap:esquemas} y, en
formato legible por máquina, en

\vspace{2pt}
\noindent\hspace*{1em}\ruta{04\_Docs/auditoria/H4\_esquemas\_derivados.csv}
\vspace{2pt}

\subsection{Tipos de dato y su representación}

\noindent\begin{tabularx}{\linewidth}{l l >{\RaggedRight\arraybackslash}X}
\toprule
\textbf{Tipo} & \textbf{En Parquet} & \textbf{Convención} \\
\midrule
Entero & \fn{int32} & Códigos DANE, conteos y años. Nunca se guardan como texto ni con decimales. \\
Decimal & \fn{double} & Tasas, índices, áreas y montos. Sin redondear en el derivado: el redondeo es cosa de la presentación. \\
Texto & \fn{string} & Nombres, categorías y etiquetas. En UTF-8 y en forma normalizada NFC. \\
Lógico & \fn{bool} & Marcas internas. No se publican. \\
\bottomrule
\end{tabularx}

\vspace{6pt}
Las columnas etiquetadas de Stata (\fn{haven\_labelled}) se guardan como su tipo
base y la etiqueta se pierde: ningún paso del flujo la usa, y conservarla
obligaría a un formato que solo Stata entiende.

\begin{regla}[title={Regla 9 — el derivado no redondea}]
Un derivado guarda el valor con toda su precisión. El redondeo ocurre una sola
vez, al escribir la hoja de Excel o la etiqueta de una figura, y queda
registrado en el formato de la celda. Redondear antes propaga el error a todos
los agregados que se calculen después.
\end{regla}

\subsection{Valores ausentes, ceros y códigos centinela}

\begin{itemize}[leftmargin=*, itemsep=2pt]
  \item En los derivados, un dato ausente es \fn{NA}. No se usan centinelas
    numéricos (\fn{-99}, \fn{999}) ni cadenas vacías para representarlo.
  \item En las hojas de Excel, un \fn{NA} se escribe como \textbf{celda vacía}.
  \item En los mapas, un municipio sin dato se pinta en gris \fn{\#E8E8E6} y la
    leyenda lo declara.
  \item Un \textbf{cero real} se escribe como \fn{0} y significa que el fenómeno
    se midió y no existe. No es lo mismo que la ausencia de medición.
\end{itemize}

\begin{advertencia}[title={Dónde el cero y el vacío se confunden con facilidad}]
En coca, EVOA, centrales de generación y áreas protegidas, la fuente solo lista
los municipios donde el fenómeno existe. La ausencia de fila significa cero. En
catastro, gobierno digital y restitución de tierras significa lo contrario: la
fuente no cubre ese municipio. La tabla de la sección~\ref{sec:trazabilidad}
separa los dos casos.
\end{advertencia}

\subsection{Unidades, escalas y precisión}

\noindent\begin{tabularx}{\linewidth}{>{\RaggedRight\arraybackslash}p{4.3cm} >{\RaggedRight\arraybackslash}X}
\toprule
\textbf{Familia} & \textbf{Unidad y escala} \\
\midrule
Áreas territoriales & km². Las áreas protegidas también en km², aunque el Anexo~1 las publicó en hectáreas. \\
Población & Personas. Proyecciones a 2025 salvo indicación en contrario. \\
Valor agregado & Miles de millones de pesos constantes de 2015; el per cápita, en pesos. \\
Avalúos catastrales & Pesos corrientes. \\
Proporciones y coberturas & Fracción en $[0,1]$ en el archivo; formato de porcentaje en la celda. \\
Tasas epidemiológicas & Por 100.000 habitantes; la mortalidad infantil por 1.000 nacidos vivos; la leishmaniasis sobre población \emph{rural}. \\
Índices (MDM, IDF, ICM, IMCA, IGD) & Escala 0--100. No se convierten a fracción. \\
Vías & Km por km² (IDV) y km por 1.000 habitantes (VPC). \\
Producción agrícola & Toneladas; áreas en hectáreas; rendimiento en t/ha. \\
Coca y oro de aluvión & Hectáreas. \\
Capacidad de generación & Megavatios (MW). \\
\bottomrule
\end{tabularx}

\begin{advertencia}[title={Una unidad que hay que declarar siempre}]
Un indicador que llega en escala 0--100 y se publica como fracción, o al revés,
produce cifras cien veces mayores o menores sin que nada falle. El flujo
convierte explícitamente en cada caso y lo anota en el script; al citar una
cifra, conviene mirar el formato de la celda antes que el número.
\end{advertencia}

\subsection{Referencia espacial}

\begin{itemize}[leftmargin=*, itemsep=2pt]
  \item \textbf{Cartografía base:} DANE, Marco Geoestadístico Nacional
    (\ruta{01\_Data/\allowbreak 00\_Inputs/\allowbreak maps/\allowbreak MGN\_MPIO\_POLITICO.shp}), 125 municipios de
    Antioquia.
  \item \textbf{Sistema de referencia:} WGS~84 geográfico, EPSG:4326.
  \item \textbf{Llave de cruce:} \fn{MPIO\_CCDGO} del shapefile contra
    \fn{ind\_mpio}. Nunca por nombre.
  \item \textbf{Codificación:} los shapefiles no traen \fn{.cpg}, así que el
    encoding se declara explícitamente al leerlos; sin eso los nombres pierden
    las tildes.
  \item \textbf{Etiquetas:} se ubican con \fn{st\_point\_on\_surface} y no con el
    centroide, que en un municipio en herradura puede caer sobre el vecino.
\end{itemize}

\subsection{Formatos, codificación y ordenación}

\noindent\begin{tabularx}{\linewidth}{l >{\RaggedRight\arraybackslash}X}
\toprule
\textbf{Elemento} & \textbf{Convención} \\
\midrule
Derivados & Parquet con compresión \fn{zstd}. Un solo formato, columnar y tipado, legible desde R, Python y Stata~18+. \\
Salidas tabulares & \fn{.xlsx}, una hoja por tabla, encabezado en la fila~1, panel congelado. \\
Figuras y mapas & PNG a 300\,dpi y PDF vectorial, ambos al tamaño final de impresión. \\
Evidencia de auditoría & CSV en UTF-8, comparables entre corridas con un \fn{diff}. \\
Codificación de texto & UTF-8 en todo el proyecto. La verificación estática rechaza nombres de archivo en forma descompuesta (NFD), que rompen en Linux. \\
Ordenación alfabética & Sobre el texto sin tildes, para que no dependa del \emph{locale}: en \fn{C}, «Gómez Plata» caería después de «Guadalupe». \\
Nombres de archivo & Los de las fuentes se respetan tal cual, con espacios y tildes. Lo que crea el pipeline no lleva ni espacios ni tildes. \\
\bottomrule
\end{tabularx}

\subsection{Nomenclatura}

\noindent\begin{tabularx}{\linewidth}{l >{\RaggedRight\arraybackslash}X}
\toprule
\textbf{Objeto} & \textbf{Patrón} \\
\midrule
Script de sección & \fn{NN\_nombre.R} en \ruta{03\_tablas} y \ruta{04\_figuras} \\
Función de sección & \fn{tabla\_nombre(prov)} y \fn{figuras\_nombre(prov, tabla)} \\
Carpeta de salida & \ruta{03\_Outputs/<Provincia>/<NN\_Sección>/} \\
Figura & \fn{fig\_NN\_slug.\{png,pdf\}}, slug en minúsculas sin tildes \\
Mapa & \fn{mapa\_NN\_slug.\{png,pdf\}}, donde \fn{NN} es la sección de destino \\
Hoja de datos de figura & \fn{\_fig\_NN} para figuras y \fn{\_mapaNN\_slug} para mapas \\
Columna del código & \fn{minúsculas\_con\_guion\_bajo}; el encabezado legible se pone al escribir la hoja \\
\bottomrule
\end{tabularx}

\subsection{Versionado, acceso y confidencialidad}

\begin{itemize}[leftmargin=*, itemsep=2pt]
  \item \textbf{Qué se versiona:} el código, las fuentes que caben, los
    tableros curados, este anexo y su evidencia. \textbf{Qué no:} las salidas
    del pipeline —se regeneran en minutos— y los archivos demasiado grandes o
    sensibles.
  \item \textbf{Huella:} cada corrida de la auditoría registra tamaño, fecha y
    MD5 de cada insumo. Es lo que permite afirmar, meses después, que una cifra
    salió exactamente de esos bytes.
  \item \textbf{Microdatos:} los de la ECV~2023 y la encuesta de percepción se
    procesan localmente; solo se versiona el derivado agregado. El de percepción
    conserva ocho variables de las 1.363 del original.
  \item \textbf{Nivel de publicación:} todo lo que el flujo publica es agregado
    municipal o superior. Ninguna salida contiene registros individuales.
\end{itemize}

% =============================================================================
\section{Protocolos de uso de datos}
% =============================================================================

\subsection{Reglas de manipulación}

\begin{regla}[title={Regla 2 — las fuentes no se tocan}]
Ningún archivo de \ruta{00\_Inputs} se edita, se limpia ni se renombra a mano.
Toda transformación ocurre en un script de \ruta{02\_Code/01\_homogeneizacion},
para que quede registrada, revisable y repetible. Un archivo corregido a mano
es una corrección que nadie puede auditar y que se pierde en la siguiente
entrega de la fuente.
\end{regla}

\begin{regla}[title={Regla 3 — los nombres de archivo son un contrato}]
Muchas fuentes traen espacios, tildes y mayúsculas inconsistentes
(\fn{POBLACION MUNICIPAL.xlsx}, \fn{códigos\_municipios\_clean.dta},
\fn{Pérdida de cobertura árborea.xlsx}). No se renombran: son el contrato entre
el código y las entregas del equipo. Un derivado conserva incluso el error
tipográfico «DEFICT» porque es el archivo vigente.
\end{regla}

\subsection{Datos sensibles y no versionados}

Dos clases de archivo no entran al repositorio:

\begin{itemize}[leftmargin=*, itemsep=2pt]
  \item \textbf{Microdatos y encuestas anonimizadas} —los microdatos de la ECV
    2023 y la encuesta de percepción de seguridad 2018--2025—, que se procesan
    localmente y de los que solo se versiona el derivado agregado. El derivado
    de percepción, por ejemplo, contiene únicamente las ocho variables que el
    análisis usa, no las 1\,363 del archivo original.
  \item \textbf{Fuentes demasiado grandes} para el control de versiones. El
    flujo se detiene con un mensaje que dice exactamente cuál falta y continúa
    con el resto.
\end{itemize}

\begin{regla}[title={Regla 4 — la encuesta de percepción no es provincial}]
La encuesta de percepción de seguridad es representativa a nivel de
\emph{subregión}, no de provincia ni de municipio. Sus resultados se publican
rotulados por subregión y deben leerse así en el informe. Presentarlos como
cifras provinciales sería atribuir a un territorio una medición que no lo
representa.
\end{regla}

\subsection{Huella de las fuentes}

Cada corrida de la auditoría registra el tamaño, la fecha de modificación y la
huella MD5 de cada insumo en
\ruta{04\_Docs/auditoria/B1\_inventario\_insumos.csv}. Esa huella es lo que
permite afirmar, meses después, que una cifra publicada salió exactamente de
esos bytes. Sin ella, «actualizamos la fuente» es una afirmación que nadie puede
verificar.

% =============================================================================
\section{Limpieza y homogeneización}
% =============================================================================

\subsection{Identificación territorial}

Todo cruce entre fuentes se hace por \textbf{código DANE}, nunca por nombre de
municipio. Cuando la fuente no trae código —las plantas de generación de XM y el
catastro minero de la ANM son los dos casos— el nombre se normaliza (mayúsculas,
sin tildes, espacios compactados, con las equivalencias conocidas del tipo
\emph{San Pedro} $\to$ \emph{San Pedro de los Milagros}) y se cruza contra el
listado oficial.

El proyecto distingue tres formas del nombre de un municipio y no las mezcla:

\begin{itemize}[leftmargin=*, itemsep=2pt]
  \item \fn{nvl\_label} — forma normalizada, en mayúsculas y sin tildes, para
    cruzar entre fuentes.
  \item \fn{municipio} — nombre de presentación, con tildes y preposiciones en
    minúscula («Carolina del Príncipe»), para tablas y figuras.
  \item El nombre tal como venga en cada fuente, que no se usa para nada más que
    para el cruce.
\end{itemize}

\begin{advertencia}[title={Por qué no se deriva el nombre de presentación}]
Aplicar una función de capitalización sobre el nombre normalizado produce
«Carolina Del Principe». El nombre de presentación se toma del listado oficial,
no se calcula.
\end{advertencia}

\subsection{Lectura de encabezados}

Los encabezados de las fuentes cambian entre entregas y entre herramientas: la
misma columna aparece como \fn{AREA KM2}, \fn{AREAKM2} o \fn{area\_km2} según
quién lea el archivo. El flujo no fija el nombre exacto: lo busca de forma
tolerante a acentos, mayúsculas, espacios y separadores, y aborta con un mensaje
que enumera las columnas disponibles cuando ninguna coincide.

Cuando el encabezado es irrecuperable —celdas combinadas, encabezados repetidos
por vigencia, títulos en dos filas— la columna se toma \emph{por posición} y esa
decisión queda anotada en el script. Es el caso del catastro, del valor agregado
municipal, del aseguramiento al SGSSS, de la natalidad y de la mortalidad
infantil.

\subsection{Conversión numérica}

Las fuentes mezclan convenciones: separador decimal de coma (\fn{1.234,5}) y de
punto (\fn{1,234.5}), notación científica (\fn{5.37E-2}) y números con símbolos.
El flujo decide el separador decimal por contexto, en este orden:

\begin{enumerate}[leftmargin=*, itemsep=1pt]
  \item Un separador repetido siempre es de miles (\fn{1.234.567}).
  \item Si aparecen los dos, el decimal es el último.
  \item Un punto solo es de miles cuando va seguido de exactamente tres dígitos
    (\fn{7.550} son siete mil quinientos cincuenta, no 7,55).
  \item Una coma sola siempre es decimal.
  \item La notación científica se deja pasar intacta.
\end{enumerate}

\begin{regla}[title={Regla 5 — cero y vacío no son lo mismo}]
Un municipio sin dato queda vacío, no en cero. En los mapas va en gris y la
leyenda lo declara. Confundirlos convierte la ausencia de medición en una
afirmación sobre el territorio.
\end{regla}

% =============================================================================
\section{Reglas de agregación}
% =============================================================================

Cada tabla provincial añade, según el indicador, una fila de provincia, una de
subregión y una de departamento. La fórmula no es la misma para todos, y la
elección no es cosmética: cambia el número publicado.

\subsection{El principio}

\begin{regla}[title={Regla 6 — el denominador del indicador manda}]
Un indicador medido \emph{por habitante} se pondera por población; uno medido
\emph{por nacimiento} —bajo peso al nacer, mortalidad infantil— se pondera por
nacimientos; uno medido \emph{por área} se pondera por área; uno escolar, por
población en edad escolar. Un conteo se suma. Una razón entre dos conteos se
recalcula sobre las sumas, no se promedia.
\end{regla}

La consecuencia práctica: donde el informe publicado promedió tasas municipales,
el flujo calcula la tasa agregada, que da más peso a los municipios donde vive
más gente. Son cifras distintas y la diferencia está documentada.

\subsection{Indicadores que no se agregan}

Cuatro familias se publican sin fila de total, por razones de fondo:

\begin{description}[leftmargin=*, style=nextline, itemsep=3pt]
  \item[Gini] Promediar los Gini municipales ignora la desigualdad
    \emph{entre} municipios y la subestima. Solo se muestran los valores
    oficiales; donde la encuesta no publica agregado, la fila queda vacía.
  \item[Ley 617] Es un cociente institucional de cada entidad, medido contra el
    techo legal de su categoría. No existe un indicador 617 provincial, y el
    departamental corresponde a la Gobernación, no a la suma de sus municipios.
  \item[MDM e IDF] El DNP compara cada municipio dentro de su grupo de
    capacidades iniciales. Un promedio entre grupos no tiene lectura.
  \item[Uso adecuado del suelo rural] La fuente cubre \cifraMunicipios{} de los
    125 municipios y la ausencia no es aleatoria. Un promedio del subconjunto
    con dato se leería como el valor del territorio completo.
\end{description}

\subsection{Verificación de las agregaciones}

La auditoría recomputa cada fila de total provincial —\cifraAgregados{} celdas
en las once provincias— y comprueba con qué fórmula se reproduce.

\begin{table}[htbp]\centering
\caption{Con qué fórmula se reproduce cada valor agregado publicado}
\small
% Generado por 02_Code/99_checks/generar_tablas_anexo.R — no editar a mano.
\begin{tabularx}{\linewidth}{X r r}
\toprule
\textbf{Fórmula que reproduce el valor publicado} & \textbf{Celdas} & \textbf{\% del total} \\
\midrule
Suma de los municipios & 4.153 & 53,7 \\
Otro ponderador declarado en el script & 2.834 & 36,7 \\
Promedio simple & 579 & 7,5 \\
Promedio ponderado por población & 144 & 1,9 \\
Mediana & 20 & 0,3 \\
\bottomrule
\addlinespace[2pt]
\multicolumn{3}{@{}p{\linewidth}@{}}{\footnotesize Verificación sobre las filas de total PROVINCIAL de todas las hojas de las once provincias. «Otro ponderador» reúne los casos en que el pipeline pondera por nacimientos, población escolar, área o matrícula, o recalcula una tasa agregada sobre la suma de numeradores y denominadores.} \\
\end{tabularx}

\end{table}

Sobre esas \cifraAgregados{} celdas quedan \cifraBanderasRojas{} que no se
reproducen con ninguna fórmula estándar \emph{y} caen fuera del rango de sus
municipios. Las cuatro están explicadas y son deliberadas: en la tabla de coca,
la fila provincial usa como denominador el área de toda la provincia —no solo la
de los municipios con cultivo—; en la de áreas protegidas, la participación de
la provincia se calcula sobre el total departamental. Ninguna otra celda del
proyecto carece de explicación.

% =============================================================================
\section{Trazabilidad y verificación}
\label{sec:trazabilidad}
% =============================================================================

\subsection{Los cuatro controles}

\begin{description}[leftmargin=*, style=nextline, itemsep=3pt]
  \item[\fn{check\_pipeline.R}] Revisión estática: rutas absolutas, insumos
    citados que no existen, nombres de archivo en forma descompuesta o que solo
    difieren en mayúsculas, secciones sin script.
  \item[\fn{comparar\_referencia.R}] Compara las tablas generadas contra las
    cifras del informe publicado, municipio por municipio y a la precisión con
    que se publicaron.
  \item[\fn{auditoria\_end\_to\_end.R}] Procedencia, huella de las fuentes,
    cobertura municipal, recomputación de agregados, rangos plausibles y
    trazabilidad de cada figura.
  \item[\fn{diccionario\_indicadores.R}] Cruza lo que el flujo publica con lo
    que el inventario declara y con las fichas técnicas existentes.
\end{description}

\subsection{Trazabilidad de las figuras}

Cada figura y cada mapa deja sus datos exactos en una hoja auxiliar del mismo
libro de Excel de su sección, con el prefijo \fn{\_}. Cualquier cifra que
aparezca dibujada se puede leer como número y rastrear hasta su tabla y de ahí
hasta su fuente. La auditoría verifica que la hoja exista para cada figura
exportada, y que cada figura del código lleve nota de fuente y subtítulo.

\begin{table}[htbp]\centering
\caption{Figuras y mapas por sección}
\small
% Generado por 02_Code/99_checks/generar_tablas_anexo.R — no editar a mano.
\begin{tabularx}{\linewidth}{X r r r}
\toprule
\textbf{Sección} & \textbf{Figuras por provincia} & \textbf{Mapas por provincia} & \textbf{Archivos generados} \\
\midrule
01 Generalidades & 1 & 2 & 33 \\
02 Demografia & 5 & 1 & 66 \\
03 Ordenamiento & 14 & 1 & 165 \\
04 Gobernabilidad & 5 & 0 & 55 \\
05 Economia & 15 & 1 & 176 \\
06 Desarrollo Rural & 5 & 0 & 55 \\
07 Ambiental & 7 & 1 & 88 \\
08 Educacion & 3 & 1 & 44 \\
09 Salud & 4 & 1 & 55 \\
10 Seguridad & 5 & 2 & 77 \\
\bottomrule
\addlinespace[2pt]
\multicolumn{4}{@{}p{\linewidth}@{}}{\footnotesize Cada pieza se exporta en PNG a 300 dpi y en PDF vectorial, y deja sus datos exactos en una hoja del .xlsx de la sección.} \\
\end{tabularx}

\end{table}

\subsection{Cobertura municipal}

La auditoría comprueba, hoja por hoja, que los municipios listados sean
exactamente los de la provincia. Un municipio ajeno sería una cifra atribuida a
un territorio que no le corresponde: no se encontró ninguno. Municipios
faltantes sí hay, y no todos significan lo mismo.

\begin{table}[htbp]\centering
\caption{Hojas donde falta la fila de algún municipio}
\small
% Generado por 02_Code/99_checks/generar_tablas_anexo.R — no editar a mano.
\begin{tabularx}{\linewidth}{X r r}
\toprule
\textbf{Hoja} & \textbf{Provincias afectadas} & \textbf{Municipios sin fila} \\
\midrule
evoa & 11 & 75 \\
coca & 10 & 72 \\
energia\_\allowbreak{}centrales & 11 & 53 \\
detalle & 8 & 19 \\
catastro & 2 & 10 \\
extranjeros\_\allowbreak{}municipio & 1 & 1 \\
gobierno\_\allowbreak{}digital & 1 & 1 \\
srtdaf & 1 & 1 \\
tasas & 1 & 1 \\
\bottomrule
\addlinespace[2pt]
\multicolumn{3}{@{}p{\linewidth}@{}}{\footnotesize Un municipio sin fila no es un cero: es un municipio del que la fuente no dice nada. En coca, EVOA, centrales de generación y áreas protegidas la ausencia es informativa —el fenómeno no existe allí—; en catastro, gobierno digital y restitución es un vacío de la fuente.} \\
\end{tabularx}

\end{table}

% =============================================================================
\section{Los documentos generados}
\label{sec:documentos}
% =============================================================================

Además de las tablas y las figuras, el flujo produce los textos: once
diagnósticos provinciales y un documento comparativo. Se generan con el mismo
principio que el resto del proyecto —ninguna cifra se escribe a mano— y se
verifican con un control propio.

\subsection{Qué se genera}

\begin{description}[leftmargin=*, style=nextline, itemsep=3pt]
  \item[Diagnóstico provincial] \ruta{04\_Docs/borradores/<Provincia>/}. Sigue
    la estructura del Anexo~1 —capítulo, diez secciones y dieciséis
    subsecciones— con sus \cifraFiguras{} figuras y mapas repartidos, seis
    tablas y la prosa redactada. Los estilos salen del propio Anexo~1, de modo
    que un borrador se puede pegar en el documento maestro sin recolocar
    niveles.
  \item[Comparativo de las once] \ruta{04\_Docs/comparativo/}. Lee las once
    provincias juntas: cuánto del departamento cubre el sistema provincial, cómo
    se ordena cada provincia en cada dimensión y qué relaciones aparecen entre
    indicadores. Se apoya en un panel de \cifraPanelIndicadores{} magnitudes por
    provincia, extraídas de las mismas hojas que publican los informes.
\end{description}

\begin{regla}[title={Regla 10 — la prosa no transcribe, calcula}]
Ninguna cifra de estos documentos la teclea una persona. Cada número se calcula
desde la tabla publicada en el momento de escribir la frase, con el formateador
del proyecto. Corregir una cifra equivocada es corregir el flujo y regenerar,
nunca editar el documento.
\end{regla}

\subsection{Hasta dónde llega la redacción generada}

El texto escribe el enunciado factual —quién concentra, quién está arriba y
abajo, cómo se compara la Provincia con su subregión y con el departamento, qué
relación aparece entre dos indicadores— y se detiene donde la frase natural
exigiría conocimiento que el flujo no tiene. Ahí deja una nota
\textit{Verificar con el equipo}: por qué un municipio se comporta como se
comporta, qué corredor vial explica un patrón, qué actor está detrás de una
economía ilegal. Son veintidós por documento y señalan trabajo pendiente, no
huecos disimulados.

\subsection{Auditoría de la redacción}

\begin{regla}[title={Regla 11 — cada cifra impresa queda registrada}]
Durante la generación, los formateadores anotan cada número que imprimen. La
auditoría vuelve a generar la prosa con ese registro activo y comprueba que
todo número del texto esté en él. La pregunta deja de ser «¿existe esta cifra en
los datos?» —que solo se puede responder buscándola en un conjunto grande, con
poca potencia— y pasa a ser «¿la calculó el flujo?», que se responde exacto.
\end{regla}

\begin{table}[htbp]\centering
\caption{Resultado de la auditoría de la redacción}
\small
\noindent\begin{tabularx}{\linewidth}{X r}
\toprule
\textbf{Comprobación} & \textbf{Resultado} \\
\midrule
Cifras auditadas en la prosa de los once documentos & \cifraProsaAuditada{} \\
Cifras sin procedencia registrada & \cifraProsaSinOrigen{} \\
Proporción con procedencia & \cifraProsaPct{}\,\% \\
\addlinespace
Cifras falsas inyectadas en el control & \cifraFalsasProbadas{} \\
Poder de detección medido & \cifraPoderDeteccion{}\,\% \\
\bottomrule
\end{tabularx}
\end{table}

\begin{hallazgo}[title={Una comprobación que nunca falla no demuestra nada}]
La auditoría mide su propio poder: inyecta cifras falsas —valores reales
alterados un 3,7\,\%, plausibles en forma y magnitud— y reporta qué proporción
detecta. Sin esa medida, un resultado del 100\,\% podría significar que todo
está bien o que la comprobación no discrimina.

La primera versión de esta auditoría buscaba cada cifra en el conjunto de
valores publicados y alcanzaba apenas un 46\,\% de poder: con miles de valores
admitidos, casi cualquier número coincidía por azar. El registro de procedencia
lo llevó a \cifraPoderDeteccion{}\,\%. El residuo corresponde a cifras alteradas
que, al redondearse, coinciden con otra que el flujo sí imprimió.
\end{hallazgo}

\subsection{Cómo se generan}

\begin{verbatim}
Rscript 02_Code/06_comparativo/generar_comparativo.R   # panel + figuras + docx
Rscript 02_Code/05_documento/generar_borrador.R        # los once diagnósticos
Rscript 02_Code/99_checks/auditoria_redaccion.R        # auditoría + control
\end{verbatim}

% =============================================================================
\section{Estrategia de actualización}
% =============================================================================

\subsection{Actualizar una fuente}

\begin{enumerate}[leftmargin=*, itemsep=3pt]
  \item Reemplazar el archivo en \ruta{01\_Data/00\_Inputs} \textbf{conservando
    el nombre}. Si el nombre cambia, hay que actualizar el script que lo lee.
  \item Correr \fn{Rscript 02\_Code/99\_checks/check\_pipeline.R}. Detecta al
    instante un insumo que dejó de existir o un encabezado que se movió.
  \item Correr el flujo completo: \fn{Rscript 02\_Code/run\_all.R}.
  \item Correr la auditoría y comparar contra la corrida anterior: los CSV de
    \ruta{04\_Docs/auditoria/} son diffables y muestran qué cambió.
  \item Regenerar las tablas de este anexo y recompilarlo.
  \item Registrar en el historial del repositorio qué fuente se actualizó, a qué
    corte y qué cifras se movieron.
\end{enumerate}

\begin{regla}[title={Regla 7 — actualizar es cambiar de corte, y hay que decirlo}]
Cambiar una fuente cambia las cifras publicadas. El informe debe declarar el
corte de cada indicador; dos indicadores de secciones distintas pueden referirse
a años distintos y el lector tiene que poder saberlo.
\end{regla}

\subsection{Cortes que ya no coinciden con el informe}

Varios indicadores usan hoy cortes posteriores a los publicados en 2026: las
cifras de salud (bajo peso, enfermedades transmitidas por vectores, suicidios,
aseguramiento) usan cortes de 2023 y de diciembre de 2025; el valor agregado
está a precios constantes de 2015 mientras el informe salió a corrientes; el uso
del suelo rural y las víctimas del conflicto vienen de versiones posteriores de
la fuente. El flujo entrega la cifra actual, que no siempre es la del informe.

\subsection{Qué hay que rehacer a mano}

Los \cifraCurados{} insumos curados no se regeneran con el flujo. Actualizarlos
exige rehacer el tablero en Excel a partir de la fuente original. Mientras eso
siga así, una actualización del diagnóstico no es una operación de un solo
comando.

% =============================================================================
\section{Fichas técnicas y diccionario de indicadores}
% =============================================================================

\subsection{Estado actual}

El repositorio conserva 65 fichas técnicas de indicador (más un juego de fichas
municipales POTA, que son documentos de contexto y no fichas de indicador). Cada
ficha define el indicador, su fórmula, su unidad, su escala de variación, el
sentido de su interpretación y su nivel de desagregación: es exactamente el
formato que este proyecto necesita.

El problema no es su calidad sino su cobertura.

\begin{hallazgo}[title={Las fichas describen otro sistema de indicadores}]
El flujo publica \cifraIndicadores{} indicadores distintos. Solo
\cifraConFicha{} tienen una ficha técnica que corresponda por nombre. Las fichas
existentes describen la familia de indicadores de la Encuesta de Calidad de Vida
y del Anuario Estadístico —calidad de vida, demografía, fuerza laboral,
educación—, mientras que el diagnóstico reporta además MDM, IDF, ICM, IMCA,
IRCA, IMRC, EVOA, IRV, gobierno digital, catastro, vías, internet fijo, uso del
suelo, títulos mineros, capacidad de generación, hectáreas de coca, hechos
victimizantes y restitución de tierras, ninguno de los cuales tiene ficha.
\end{hallazgo}

\begin{table}[htbp]\centering
\caption{Cobertura documental de los indicadores publicados}
\small
% Generado por 02_Code/99_checks/generar_tablas_anexo.R — no editar a mano.
\begin{tabularx}{\linewidth}{X r r r}
\toprule
\textbf{Sección} & \textbf{Indicadores publicados} & \textbf{Con ficha técnica} & \textbf{En el inventario} \\
\midrule
01 Generalidades & 2 & 0 & 0 \\
02 Demografia & 44 & 0 & 2 \\
03 Ordenamiento & 57 & 2 & 0 \\
04 Gobernabilidad & 18 & 0 & 0 \\
05 Economia & 87 & 9 & 2 \\
06 Desarrollo Rural & 30 & 0 & 0 \\
07 Ambiental & 23 & 0 & 0 \\
08 Educacion & 3 & 1 & 0 \\
09 Salud & 13 & 0 & 0 \\
10 Seguridad & 53 & 0 & 0 \\
\bottomrule
\addlinespace[2pt]
\multicolumn{4}{@{}p{\linewidth}@{}}{\footnotesize El cruce es EXACTO sobre el nombre normalizado del indicador. Un indicador que no case queda sin respaldo documental, que es lo que este anexo pide corregir.} \\
\end{tabularx}

\end{table}

\subsection{Qué hacer con eso}

El inventario \fn{INVENTARIO\_VARIABLES\_v\_2.xlsx} ya tiene la estructura
correcta —fuente, periodo, periodicidad, sentido, unidad y estado de
actualización por indicador— pero tampoco cubre lo que el flujo publica. La ruta
más corta es completarlo con los indicadores realmente publicados, usando como
llave el encabezado exacto de la tabla, y generar las fichas faltantes desde ahí
en vez de escribirlas una a una.

El diccionario completo de lo que hoy se publica —sección, hoja, encabezado, en
cuántas provincias aparece y si tiene ficha— está en
\ruta{04\_Docs/auditoria/H1\_diccionario\_indicadores.csv}.

% =============================================================================
\section{Limitaciones conocidas}
\label{sec:limitaciones}
% =============================================================================

Esta sección existe para que nadie descubra estos límites después de publicar.

\subsection{Insumos que se arman a mano}

\begin{hallazgo}[title={El mayor riesgo de reproducibilidad}]
\cifraCurados{} insumos se construyen a mano en Excel y ningún script los
regenera. Alimentan secciones completas: gobernabilidad (MDM, IDF, Ley 617,
ICM), ambiental (áreas protegidas, IRCA, emergencias, IMRC), desarrollo rural
(inventario pecuario, cultivos), seguridad (coca, EVOA, hechos victimizantes,
restitución), economía (valor agregado, visitantes extranjeros) y educación. Si
uno se pierde, hay que rehacerlo a mano y no hay forma de verificar que quedó
igual.
\end{hallazgo}

\subsection{Derivados cuyo productor no es ejecutable}

\begin{hallazgo}[title={Cuatro cadenas que no se pueden rehacer}]
Los scripts de \ruta{02\_Code/01\_homogeneizacion/externos/} documentan cómo se
construyeron cuatro insumos curados —vías (IDV y VPC), suicidios e intentos,
índice de envejecimiento e IRCA—, pero leen rutas del computador de quien los
escribió y no corren en este repositorio. Dos de ellos, además, no contienen
ninguna instrucción de escritura: su salida está comentada o simplemente no
existe. Hay un eslabón manual adicional sin documentar —añadir a mano una hoja
\fn{Data} en Excel al resultado— que explica el sufijo \fn{\_dicc} de esos
archivos.

Consecuencia concreta: las tasas de intento de suicidio y de suicidio consumado,
los kilómetros de vía por área y por habitante, y el índice de envejecimiento se
publican sin que exista en el repositorio código capaz de reproducirlos desde la
fuente.
\end{hallazgo}

\subsection{Validación externa limitada}

\begin{hallazgo}[title={Solo una de once provincias tiene contraste externo}]
El único informe publicado disponible como referencia es el de la provincia
Bioenergética del Norte. El comparador valida 43 columnas contra él y las 43
reproducen el informe. Para las otras diez provincias no hay contraste externo:
su corrección descansa en que el código es el mismo y en los controles internos
de este anexo. No es lo mismo y conviene decirlo así.
\end{hallazgo}

\subsection{Un indicador que la fuente publica mal}

\begin{hallazgo}[title={Cobertura neta por encima del 100\,\%}]
La cobertura neta educativa supera el 100\,\% en 284 de los 1\,750 registros
municipio-año de la fuente, con un máximo de 186,97\,\%. Una cobertura
\emph{neta} —matriculados en la edad teórica sobre población de esa edad— no
puede exceder el 100\,\% por definición. El valor viene así de la fuente: o la
serie corresponde en realidad a cobertura bruta, o el numerador de matrícula y
el denominador de proyección poblacional no son consistentes.

El flujo publica el dato sin corregirlo —corregirlo sería inventar— pero las
figuras y los mapas de esa sección llevan una nota que lo advierte. La solución
de fondo es verificar con el Ministerio de Educación Nacional qué serie es.
\end{hallazgo}

\subsection{Vacíos de cobertura}

\cifraHojasConVacios{} hojas tienen al menos un municipio sin fila. En coca,
EVOA, centrales de generación y áreas protegidas la ausencia es informativa: el
fenómeno no existe en ese municipio. En catastro, gobierno digital y restitución
de tierras es un vacío de la fuente. El caso más grande es el catastro del Área
Metropolitana, que la fuente no cubre.

\subsection{Documentación desactualizada}

Dos piezas del repositorio contradicen la implementación vigente y conviene
retirarlas o marcarlas: \ruta{02\_Code/04\_figuras/CATALOGO\_NOTAS.md} describe
una decisión ya superada —gráficos nativos de Excel, sin ggplot2— y el catálogo
\fn{catalogo\_figuras.xlsx} marca 78 figuras como «hecho» cuando el código R
genera 56 por provincia más 10 mapas. Las cabeceras de varios scripts de tablas
declaran rutas en \ruta{01\_Data/01\_Derived} para insumos que hoy viven en
\ruta{01\_Data/00\_Inputs/curados}.

% =============================================================================
\section{Cómo reproducir todo}
% =============================================================================

\begin{verbatim}
# 1. Dependencias (la primera vez)
Rscript -e 'install.packages("renv"); renv::restore()'
Rscript -e 'install.packages(c("sf", "ggrepel"))'   # cartografía

# 2. Flujo completo: fuentes -> derivados -> tablas -> figuras -> mapas
Rscript 02_Code/run_all.R

# 3. Verificación
Rscript 02_Code/99_checks/check_pipeline.R
Rscript 02_Code/99_checks/comparar_referencia.R
Rscript 02_Code/99_checks/auditoria_end_to_end.R
Rscript 02_Code/99_checks/diccionario_indicadores.R

# 4. Este anexo
Rscript 02_Code/99_checks/generar_tablas_anexo.R
cd 04_Docs && latexmk -pdf anexo_metodologico.tex
\end{verbatim}

\vspace{1em}
\noindent
\textcolor{tintatres}{\small
Todos los CSV que respaldan las cifras de este documento están en
\ruta{04\_Docs/auditoria/}. Cada uno se regenera con el comando de la
sección~\ref{sec:trazabilidad} y puede compararse entre corridas.}

% =============================================================================
\appendix
\clearpage
\section{Esquema de las tablas derivadas}
\label{ap:esquemas}
% =============================================================================

Las \cifraColumnasDerivadas{} columnas de las \cifraTablasDerivadas{} tablas
derivadas, leídas de los propios archivos Parquet. «Distintos» es el número de
valores diferentes que toma la columna —un indicador rápido de si es una llave,
una categoría o una medida— y «Ejemplo» es un valor real, que dice más sobre la
unidad y la escala que cualquier descripción.

La misma información, en formato legible por máquina, está en

\vspace{2pt}
\noindent\hspace*{1em}\ruta{04\_Docs/auditoria/H4\_esquemas\_derivados.csv}
\vspace{2pt}

\vspace{6pt}
{\scriptsize
% Generado por 02_Code/99_checks/generar_tablas_anexo.R — no editar a mano.
\begin{longtable}{>{\raggedright\arraybackslash}p{3.1cm} >{\raggedright\arraybackslash}p{3.6cm} l r r >{\raggedright\arraybackslash}p{2.7cm}}
\toprule \textbf{Tabla} & \textbf{Columna} & \textbf{Tipo} & \textbf{Distintos} & \textbf{\% vacíos} & \textbf{Ejemplo} \\ \midrule
\endfirsthead
\multicolumn{6}{@{}l}{\footnotesize\itshape (continúa de la página anterior)} \\
\toprule \textbf{Tabla} & \textbf{Columna} & \textbf{Tipo} & \textbf{Distintos} & \textbf{\% vacíos} & \textbf{Ejemplo} \\ \midrule
\endhead
\midrule \multicolumn{6}{r@{}}{\footnotesize\itshape continúa en la página siguiente} \\
\endfoot
\bottomrule
\endlastfoot
20250506 EDUCACION & ind\_\allowbreak{}mpio & decimal & 125 & 0 &    5001 \\
 & nvl\_\allowbreak{}label & texto & 125 & 0 & Medellín \\
 & clas\_\allowbreak{}cole & decimal & 43 & 0 & 42.0588 \\
 & ninis & decimal & 125 & 0 & 21.4838 \\
 & tcb\_\allowbreak{}edsup & decimal & 52 & 0 & 115.969 \\
 & tcb\_\allowbreak{}media & decimal & 122 & 0 &   98.57 \\
 & tcb\_\allowbreak{}primaria & decimal & 124 & 0 &  104.44 \\
 & tcb\_\allowbreak{}secundaria & decimal & 123 & 0 &  110.28 \\
 & tcb\_\allowbreak{}transicion & decimal & 125 & 0 &   88.55 \\
 & tdes\_\allowbreak{}media & decimal & 113 & 0 &    3.93 \\
 & tdes\_\allowbreak{}primaria & decimal & 116 & 0 &    4.85 \\
 & tdes\_\allowbreak{}secundaria & decimal & 112 & 0 &    6.75 \\
 & tdes\_\allowbreak{}transicion & decimal & 106 & 0 &    4.84 \\
 & trep\_\allowbreak{}media & decimal & 104 & 0 &    4.82 \\
 & trep\_\allowbreak{}primaria & decimal & 115 & 0 &    7.86 \\
 & trep\_\allowbreak{}secundaria & decimal & 121 & 0 &   13.73 \\
 & trep\_\allowbreak{}transicion & decimal & 85 & 0 &    1.89 \\
 & tti\_\allowbreak{}edsup & decimal & 124 & 0 & 49.3626 \\
20260504\_\allowbreak{}SEGURIDAD\_\allowbreak{}SALUD\_\allowbreak{}DEFICT\_\allowbreak{}VIVIENDA & ind\_\allowbreak{}mpio & decimal & 125 & 0 &    5001 \\
 & nvl\_\allowbreak{}label & texto & 125 & 0 & Medellín \\
 & tasa\_\allowbreak{}natalidad & decimal & 125 & 0 & 7.80153 \\
 & tasa\_\allowbreak{}mortalidad & decimal & 125 & 0 & 6.47471 \\
 & crec\_\allowbreak{}vegetativo & decimal & 125 & 0 & 1.32682 \\
 & percepción\_\allowbreak{}seguridad & decimal & 125 & 0 & 0.606802 \\
 & tasa\_\allowbreak{}delitos & decimal & 108 & 0 & 0.0677347 \\
 & IRV & decimal & 125 & 0 & 0.183561 \\
 & VBG & decimal & 125 & 0 & 426.041 \\
 & VBG\_\allowbreak{}jovenes & decimal & 125 & 0 & 557.056 \\
 & VBG\_\allowbreak{}adolescentes & decimal & 124 & 0 & 708.421 \\
 & homicidios & decimal & 124 & 0 & 13.0339 \\
 & vinculacion\_\allowbreak{}nna & decimal & 26 & 0 &       0 \\
 & mortalidad\_\allowbreak{}desnutrición & decimal & 17 & 0 & 1.67496 \\
 & bajo\_\allowbreak{}peso & decimal & 114 & 0 & 12.3695 \\
 & desnutrición\_\allowbreak{}aguda & decimal & 115 & 0 & 562.471 \\
 & controles\_\allowbreak{}prenatales & decimal & 111 & 0 & 4.00774 \\
 & fecundidad\_\allowbreak{}niñas & decimal & 114 & 0 & 1.18822 \\
 & fecundidad\_\allowbreak{}adolescentes & decimal & 125 & 0 & 27.7947 \\
 & camas & decimal & 125 & 0 &  283.69 \\
 & IRA & decimal & 42 & 0 & 6.09654 \\
 & dengue & decimal & 90 & 0 & 12.3129 \\
 & malaria & decimal & 60 & 0 & 22.1426 \\
 & leishmaniasis & decimal & 89 & 0 & 21.4145 \\
 & acueducto\_\allowbreak{}rural & decimal & 116 & 0 & 0.920245 \\
 & alcantarillado\_\allowbreak{}rural & decimal & 118 & 0 & 4.29448 \\
 & acueducto\_\allowbreak{}urbano & decimal & 92 & 0 & 1.24469 \\
 & alcantarillado\_\allowbreak{}urbano & decimal & 105 & 0 & 1.56685 \\
 & deficit\_\allowbreak{}cuali & decimal & 125 & 0.8 & 13.6663 \\
 & deficit\_\allowbreak{}cuanti & decimal & 124 & 0.8 &  1.4732 \\
 & hacinamiento\_\allowbreak{}rural & decimal & 57 & 0 & 0.306748 \\
 & hacinamiento\_\allowbreak{}urbano & decimal & 47 & 0 & 0.644311 \\
 & Población\_\allowbreak{}2023 & decimal & 125 & 0 & 2.5953e+06 \\
 & IRCC & decimal & 125 & 0 & 0.214702 \\
 & afectados\_\allowbreak{}dn & decimal & 123 & 0 & 289.788 \\
 & perdida\_\allowbreak{}ca & decimal & 70 & 0 &     0.1 \\
codigos\_\allowbreak{}provincias & ind\_\allowbreak{}mpio & entero & 88 & 0 &    5001 \\
 & nvl\_\allowbreak{}label & texto & 88 & 0 & MEDELLIN \\
 & municipio & texto & 88 & 0 & Medellín \\
 & subregion & texto & 6 & 0 & VALLE DE ABURRA \\
 & provincia & texto & 11 & 0 & AREA METROPOLITANA \\
 & id\_\allowbreak{}provincia & entero & 11 & 0 &      11 \\
ECV\_\allowbreak{}nbi\_\allowbreak{}pobreza & ind\_\allowbreak{}mpio & entero & 125 & 0 &    5001 \\
 & NomMunicipio & texto & 125 & 0 & Medellín \\
 & tot\_\allowbreak{}emp\_\allowbreak{}informal & decimal & 125 & 0 & 41.1278 \\
 & urb\_\allowbreak{}emp\_\allowbreak{}informal & decimal & 125 & 0 & 40.8952 \\
 & rur\_\allowbreak{}emp\_\allowbreak{}informal & decimal & 125 & 0 & 42.7051 \\
 & tot\_\allowbreak{}emp\_\allowbreak{}informal\_\allowbreak{}hombres & decimal & 125 & 0 &  40.835 \\
 & urb\_\allowbreak{}emp\_\allowbreak{}informal\_\allowbreak{}hombres & decimal & 125 & 0 &  40.711 \\
 & rur\_\allowbreak{}emp\_\allowbreak{}informal\_\allowbreak{}hombres & decimal & 125 & 0 &  41.662 \\
 & tot\_\allowbreak{}emp\_\allowbreak{}informal\_\allowbreak{}mujeres & decimal & 125 & 0 & 41.4949 \\
 & urb\_\allowbreak{}emp\_\allowbreak{}informal\_\allowbreak{}mujeres & decimal & 125 & 0 &  41.125 \\
 & rur\_\allowbreak{}emp\_\allowbreak{}informal\_\allowbreak{}mujeres & decimal & 125 & 0 & 44.0561 \\
 & tot\_\allowbreak{}gini\_\allowbreak{}hog & decimal & 125 & 0 & 0.49453 \\
 & urb\_\allowbreak{}gini\_\allowbreak{}hog & decimal & 125 & 0 & 0.49368 \\
 & rur\_\allowbreak{}gini\_\allowbreak{}hog & decimal & 125 & 0 & 0.491662 \\
 & tot\_\allowbreak{}gini\_\allowbreak{}lab & decimal & 125 & 0 & 0.37418 \\
 & urb\_\allowbreak{}gini\_\allowbreak{}lab & decimal & 125 & 0 & 0.372657 \\
 & rur\_\allowbreak{}gini\_\allowbreak{}lab & decimal & 125 & 0 & 0.380054 \\
 & tot\_\allowbreak{}hog\_\allowbreak{}mono\_\allowbreak{}fem & decimal & 125 & 0 & 12.7035 \\
 & urb\_\allowbreak{}hog\_\allowbreak{}mono\_\allowbreak{}fem & decimal & 110 & 0 & 12.6101 \\
 & rur\_\allowbreak{}hog\_\allowbreak{}mono\_\allowbreak{}fem & decimal & 104 & 0 &  13.369 \\
 & tot\_\allowbreak{}ins\_\allowbreak{}alim\_\allowbreak{}mod & decimal & 125 & 0 & 19.6128 \\
 & urb\_\allowbreak{}ins\_\allowbreak{}alim\_\allowbreak{}mod & decimal & 115 & 0 & 19.3633 \\
 & rur\_\allowbreak{}ins\_\allowbreak{}alim\_\allowbreak{}mod & decimal & 114 & 0 & 21.3701 \\
 & tot\_\allowbreak{}ins\_\allowbreak{}alim\_\allowbreak{}mod\_\allowbreak{}18 & decimal & 125 & 0 & 22.3153 \\
 & urb\_\allowbreak{}ins\_\allowbreak{}alim\_\allowbreak{}mod\_\allowbreak{}18 & decimal & 98 & 0 & 22.3301 \\
 & rur\_\allowbreak{}ins\_\allowbreak{}alim\_\allowbreak{}mod\_\allowbreak{}18 & decimal & 98 & 0 & 22.2222 \\
 & tot\_\allowbreak{}ins\_\allowbreak{}alim\_\allowbreak{}mod\_\allowbreak{}5 & decimal & 124 & 0 & 24.4751 \\
 & urb\_\allowbreak{}ins\_\allowbreak{}alim\_\allowbreak{}mod\_\allowbreak{}5 & decimal & 63 & 0 & 23.8987 \\
 & rur\_\allowbreak{}ins\_\allowbreak{}alim\_\allowbreak{}mod\_\allowbreak{}5 & decimal & 61 & 0 & 28.0822 \\
 & tot\_\allowbreak{}ins\_\allowbreak{}alim\_\allowbreak{}sev & decimal & 125 & 0 & 18.8734 \\
 & urb\_\allowbreak{}ins\_\allowbreak{}alim\_\allowbreak{}sev & decimal & 123 & 0 & 18.9544 \\
 & rur\_\allowbreak{}ins\_\allowbreak{}alim\_\allowbreak{}sev & decimal & 121 & 0 & 18.3027 \\
 & tot\_\allowbreak{}ins\_\allowbreak{}alim\_\allowbreak{}sev\_\allowbreak{}18 & decimal & 125 & 0 & 20.2587 \\
 & urb\_\allowbreak{}ins\_\allowbreak{}alim\_\allowbreak{}sev\_\allowbreak{}18 & decimal & 105 & 0 & 20.3883 \\
 & rur\_\allowbreak{}ins\_\allowbreak{}alim\_\allowbreak{}sev\_\allowbreak{}18 & decimal & 95 & 0 & 19.4444 \\
 & tot\_\allowbreak{}ins\_\allowbreak{}alim\_\allowbreak{}sev\_\allowbreak{}5 & decimal & 122 & 0 & 17.7431 \\
 & urb\_\allowbreak{}ins\_\allowbreak{}alim\_\allowbreak{}sev\_\allowbreak{}5 & decimal & 62 & 0 & 17.9515 \\
 & rur\_\allowbreak{}ins\_\allowbreak{}alim\_\allowbreak{}sev\_\allowbreak{}5 & decimal & 57 & 0 & 16.4384 \\
 & tot\_\allowbreak{}nini & decimal & 125 & 0 & 21.4838 \\
 & urb\_\allowbreak{}nini & decimal & 125 & 0 & 21.5816 \\
 & rur\_\allowbreak{}nini & decimal & 125 & 0 & 20.7848 \\
 & tot\_\allowbreak{}pbl\_\allowbreak{}camp & decimal & 125 & 0 & 14.7742 \\
 & urb\_\allowbreak{}pbl\_\allowbreak{}camp & decimal & 125 & 0 & 11.3785 \\
 & rur\_\allowbreak{}pbl\_\allowbreak{}camp & decimal & 116 & 0 & 37.7156 \\
 & tot\_\allowbreak{}pob\_\allowbreak{}acueducto & decimal & 123 & 0 & 98.6685 \\
 & urb\_\allowbreak{}pob\_\allowbreak{}acueducto & decimal & 93 & 0 & 98.7553 \\
 & rur\_\allowbreak{}pob\_\allowbreak{}acueducto & decimal & 117 & 0 & 98.0573 \\
 & tot\_\allowbreak{}pob\_\allowbreak{}alcantarillado & decimal & 125 & 0 & 97.9849 \\
 & urb\_\allowbreak{}pob\_\allowbreak{}alcantarillado & decimal & 103 & 0 & 99.0921 \\
 & rur\_\allowbreak{}pob\_\allowbreak{}alcantarillado & decimal & 121 & 0 &  90.184 \\
 & tot\_\allowbreak{}pob\_\allowbreak{}energia & decimal & 96 & 0 & 99.3856 \\
 & urb\_\allowbreak{}pob\_\allowbreak{}energia & decimal & 44 & 0 & 99.4435 \\
 & rur\_\allowbreak{}pob\_\allowbreak{}energia & decimal & 70 & 0 & 98.9775 \\
 & tot\_\allowbreak{}pob\_\allowbreak{}ext & decimal & 119 & 0 & 2.74881 \\
 & urb\_\allowbreak{}pob\_\allowbreak{}ext & decimal & 112 & 0 & 2.67135 \\
 & rur\_\allowbreak{}pob\_\allowbreak{}ext & decimal & 98 & 0 & 3.25773 \\
 & tot\_\allowbreak{}pob\_\allowbreak{}ext\_\allowbreak{}hombres & decimal & 116 & 0 & 2.66322 \\
 & urb\_\allowbreak{}pob\_\allowbreak{}ext\_\allowbreak{}hombres & decimal & 110 & 0 & 2.58988 \\
 & rur\_\allowbreak{}pob\_\allowbreak{}ext\_\allowbreak{}hombres & decimal & 94 & 0 &  3.1453 \\
 & tot\_\allowbreak{}pob\_\allowbreak{}ext\_\allowbreak{}mujeres & decimal & 119 & 0 & 2.82459 \\
 & urb\_\allowbreak{}pob\_\allowbreak{}ext\_\allowbreak{}mujeres & decimal & 112 & 0 &  2.7435 \\
 & rur\_\allowbreak{}pob\_\allowbreak{}ext\_\allowbreak{}mujeres & decimal & 98 & 0 & 3.35718 \\
 & tot\_\allowbreak{}pob\_\allowbreak{}gas\_\allowbreak{}natural & decimal & 112 & 0 & 79.2927 \\
 & urb\_\allowbreak{}pob\_\allowbreak{}gas\_\allowbreak{}natural & decimal & 107 & 0 & 81.5346 \\
 & rur\_\allowbreak{}pob\_\allowbreak{}gas\_\allowbreak{}natural & decimal & 97 & 0 & 63.4969 \\
 & tot\_\allowbreak{}pob\_\allowbreak{}gas\_\allowbreak{}pipeta & decimal & 125 & 0 & 18.7229 \\
 & urb\_\allowbreak{}pob\_\allowbreak{}gas\_\allowbreak{}pipeta & decimal & 122 & 0 & 16.6057 \\
 & rur\_\allowbreak{}pob\_\allowbreak{}gas\_\allowbreak{}pipeta & decimal & 121 & 0 & 33.6401 \\
 & tot\_\allowbreak{}pob\_\allowbreak{}internet & decimal & 125 & 0 & 70.4836 \\
 & urb\_\allowbreak{}pob\_\allowbreak{}internet & decimal & 122 & 0 & 71.5478 \\
 & rur\_\allowbreak{}pob\_\allowbreak{}internet & decimal & 119 & 0 & 62.9857 \\
 & tot\_\allowbreak{}pob\_\allowbreak{}ipm & decimal & 125 & 0 & 6.03684 \\
 & urb\_\allowbreak{}pob\_\allowbreak{}ipm & decimal & 122 & 0 & 6.20355 \\
 & rur\_\allowbreak{}pob\_\allowbreak{}ipm & decimal & 125 & 0 & 4.94158 \\
 & tot\_\allowbreak{}pob\_\allowbreak{}ipm\_\allowbreak{}hogares & decimal & 125 & 0 & 4.17868 \\
 & urb\_\allowbreak{}pob\_\allowbreak{}ipm\_\allowbreak{}hogares & decimal & 118 & 0 & 4.23481 \\
 & rur\_\allowbreak{}pob\_\allowbreak{}ipm\_\allowbreak{}hogares & decimal & 118 & 0 & 3.78323 \\
 & tot\_\allowbreak{}pob\_\allowbreak{}ipm\_\allowbreak{}hombres & decimal & 125 & 0 & 6.21303 \\
 & urb\_\allowbreak{}pob\_\allowbreak{}ipm\_\allowbreak{}hombres & decimal & 122 & 0 & 6.35592 \\
 & rur\_\allowbreak{}pob\_\allowbreak{}ipm\_\allowbreak{}hombres & decimal & 125 & 0 &  5.2737 \\
 & tot\_\allowbreak{}pob\_\allowbreak{}ipm\_\allowbreak{}mujeres & decimal & 125 & 0 & 5.88086 \\
 & urb\_\allowbreak{}pob\_\allowbreak{}ipm\_\allowbreak{}mujeres & decimal & 122 & 0 & 6.06862 \\
 & rur\_\allowbreak{}pob\_\allowbreak{}ipm\_\allowbreak{}mujeres & decimal & 125 & 0 & 4.64779 \\
 & tot\_\allowbreak{}pob\_\allowbreak{}mon & decimal & 125 & 0 &   24.61 \\
 & urb\_\allowbreak{}pob\_\allowbreak{}mon & decimal & 125 & 0 & 24.2448 \\
 & rur\_\allowbreak{}pob\_\allowbreak{}mon & decimal & 125 & 0 & 27.0093 \\
 & tot\_\allowbreak{}pob\_\allowbreak{}mon\_\allowbreak{}hombres & decimal & 125 & 0 &  24.151 \\
 & urb\_\allowbreak{}pob\_\allowbreak{}mon\_\allowbreak{}hombres & decimal & 125 & 0 &   23.78 \\
 & rur\_\allowbreak{}pob\_\allowbreak{}mon\_\allowbreak{}hombres & decimal & 125 & 0 & 26.5899 \\
 & tot\_\allowbreak{}pob\_\allowbreak{}mon\_\allowbreak{}mujeres & decimal & 125 & 0 & 25.0163 \\
 & urb\_\allowbreak{}pob\_\allowbreak{}mon\_\allowbreak{}mujeres & decimal & 125 & 0 & 24.6564 \\
 & rur\_\allowbreak{}pob\_\allowbreak{}mon\_\allowbreak{}mujeres & decimal & 125 & 0 & 27.3803 \\
 & tot\_\allowbreak{}pob\_\allowbreak{}nbi & decimal & 125 & 0 & 6.40851 \\
 & urb\_\allowbreak{}pob\_\allowbreak{}nbi & decimal & 125 & 0 & 6.89985 \\
 & rur\_\allowbreak{}pob\_\allowbreak{}nbi & decimal & 125 & 0 & 3.18038 \\
 & tot\_\allowbreak{}pob\_\allowbreak{}nbi\_\allowbreak{}hogares & decimal & 125 & 0 & 6.37803 \\
 & urb\_\allowbreak{}pob\_\allowbreak{}nbi\_\allowbreak{}hogares & decimal & 117 & 0 & 6.71729 \\
 & rur\_\allowbreak{}pob\_\allowbreak{}nbi\_\allowbreak{}hogares & decimal & 119 & 0 & 3.98773 \\
 & tot\_\allowbreak{}pob\_\allowbreak{}nbi\_\allowbreak{}hombres & decimal & 125 & 0 & 6.27578 \\
 & urb\_\allowbreak{}pob\_\allowbreak{}nbi\_\allowbreak{}hombres & decimal & 125 & 0 & 6.80928 \\
 & rur\_\allowbreak{}pob\_\allowbreak{}nbi\_\allowbreak{}hombres & decimal & 125 & 0 & 2.76878 \\
 & tot\_\allowbreak{}pob\_\allowbreak{}nbi\_\allowbreak{}mujeres & decimal & 125 & 0 & 6.52603 \\
 & urb\_\allowbreak{}pob\_\allowbreak{}nbi\_\allowbreak{}mujeres & decimal & 125 & 0 & 6.98005 \\
 & rur\_\allowbreak{}pob\_\allowbreak{}nbi\_\allowbreak{}mujeres & decimal & 125 & 0 & 3.54448 \\
 & tot\_\allowbreak{}pob\_\allowbreak{}recoleccion\_\allowbreak{}basuras & decimal & 125 & 0 & 97.9773 \\
 & urb\_\allowbreak{}pob\_\allowbreak{}recoleccion\_\allowbreak{}basuras & decimal & 94 & 0 & 98.1256 \\
 & rur\_\allowbreak{}pob\_\allowbreak{}recoleccion\_\allowbreak{}basuras & decimal & 112 & 0 & 96.9325 \\
 & tot\_\allowbreak{}td & decimal & 125 & 0 & 8.95271 \\
 & urb\_\allowbreak{}td & decimal & 125 & 0 & 8.77838 \\
 & rur\_\allowbreak{}td & decimal & 118 & 0 & 10.1175 \\
 & tot\_\allowbreak{}td\_\allowbreak{}15\_\allowbreak{}28 & decimal & 122 & 0 & 12.7316 \\
 & urb\_\allowbreak{}td\_\allowbreak{}15\_\allowbreak{}28 & decimal & 114 & 0 & 12.8028 \\
 & rur\_\allowbreak{}td\_\allowbreak{}15\_\allowbreak{}28 & decimal & 97 & 0 & 12.2529 \\
 & tot\_\allowbreak{}td\_\allowbreak{}hombres & decimal & 120 & 0 & 7.78851 \\
 & urb\_\allowbreak{}td\_\allowbreak{}hombres & decimal & 112 & 0 & 7.59081 \\
 & rur\_\allowbreak{}td\_\allowbreak{}hombres & decimal & 90 & 0 & 9.08562 \\
 & tot\_\allowbreak{}td\_\allowbreak{}mujeres & decimal & 123 & 0 &  10.371 \\
 & urb\_\allowbreak{}td\_\allowbreak{}mujeres & decimal & 121 & 0 & 10.2175 \\
 & rur\_\allowbreak{}td\_\allowbreak{}mujeres & decimal & 113 & 0 & 11.4196 \\
 & tot\_\allowbreak{}tgp\_\allowbreak{}mujeres & decimal & 125 & 0 & 39.9177 \\
 & urb\_\allowbreak{}tgp\_\allowbreak{}mujeres & decimal & 125 & 0 & 40.1239 \\
 & rur\_\allowbreak{}tgp\_\allowbreak{}mujeres & decimal & 125 & 0 & 38.5637 \\
 & tot\_\allowbreak{}to & decimal & 125 & 0 & 51.8484 \\
 & urb\_\allowbreak{}to & decimal & 125 & 0 & 51.8729 \\
 & rur\_\allowbreak{}to & decimal & 125 & 0 & 51.6826 \\
 & tot\_\allowbreak{}to\_\allowbreak{}hombres & decimal & 125 & 0 & 62.5465 \\
 & urb\_\allowbreak{}to\_\allowbreak{}hombres & decimal & 125 & 0 & 62.4191 \\
 & rur\_\allowbreak{}to\_\allowbreak{}hombres & decimal & 125 & 0 & 63.4098 \\
 & tot\_\allowbreak{}to\_\allowbreak{}mujeres & decimal & 125 & 0 & 42.6956 \\
 & urb\_\allowbreak{}to\_\allowbreak{}mujeres & decimal & 125 & 0 & 42.8441 \\
 & rur\_\allowbreak{}to\_\allowbreak{}mujeres & decimal & 125 & 0 &  41.695 \\
ECV\_\allowbreak{}oficial\_\allowbreak{}agregados & nivel & texto & 3 & 0 & DEPARTAMENTO \\
 & key & texto & 16 & 0 & ANTIOQUIA \\
 & terr\_\allowbreak{}label & texto & 16 & 0 & Antioquia \\
 & rur\_\allowbreak{}emp\_\allowbreak{}informal & decimal & 11 & 37.5 & 0.576232 \\
 & rur\_\allowbreak{}gini\_\allowbreak{}hog & decimal & 16 & 6.2 & 0.438508 \\
 & rur\_\allowbreak{}gini\_\allowbreak{}lab & decimal & 16 & 6.2 & 0.362559 \\
 & rur\_\allowbreak{}nini & decimal & 16 & 0 & 0.281195 \\
 & rur\_\allowbreak{}pob\_\allowbreak{}acueducto & decimal & 16 & 0 & 0.798239 \\
 & rur\_\allowbreak{}pob\_\allowbreak{}alcantarillado & decimal & 16 & 0 & 0.541232 \\
 & rur\_\allowbreak{}pob\_\allowbreak{}energia & decimal & 16 & 0 & 0.983608 \\
 & rur\_\allowbreak{}pob\_\allowbreak{}gas\_\allowbreak{}natural & decimal & 16 & 0 & 0.296645 \\
 & rur\_\allowbreak{}pob\_\allowbreak{}internet & decimal & 16 & 0 & 0.407165 \\
 & rur\_\allowbreak{}pob\_\allowbreak{}ipm & decimal & 16 & 0 & 0.177617 \\
 & rur\_\allowbreak{}pob\_\allowbreak{}ipm\_\allowbreak{}hogares & decimal & 16 & 0 & 0.153547 \\
 & rur\_\allowbreak{}pob\_\allowbreak{}nbi & decimal & 16 & 0 & 0.153621 \\
 & rur\_\allowbreak{}pob\_\allowbreak{}recoleccion\_\allowbreak{}basuras & decimal & 16 & 0 & 0.835154 \\
 & rur\_\allowbreak{}td & decimal & 11 & 37.5 & 0.064345 \\
 & rur\_\allowbreak{}td\_\allowbreak{}15\_\allowbreak{}28 & decimal & 11 & 37.5 & 0.084816 \\
 & rur\_\allowbreak{}to & decimal & 11 & 37.5 & 0.492264 \\
 & tot\_\allowbreak{}emp\_\allowbreak{}informal & decimal & 11 & 37.5 & 0.473997 \\
 & tot\_\allowbreak{}gini\_\allowbreak{}hog & decimal & 16 & 6.2 & 0.437434 \\
 & tot\_\allowbreak{}gini\_\allowbreak{}lab & decimal & 16 & 6.2 & 0.347928 \\
 & tot\_\allowbreak{}nini & decimal & 16 & 0 & 0.245564 \\
 & tot\_\allowbreak{}pob\_\allowbreak{}acueducto & decimal & 16 & 0 & 0.935796 \\
 & tot\_\allowbreak{}pob\_\allowbreak{}alcantarillado & decimal & 16 & 0 & 0.866894 \\
 & tot\_\allowbreak{}pob\_\allowbreak{}energia & decimal & 16 & 0 & 0.992587 \\
 & tot\_\allowbreak{}pob\_\allowbreak{}gas\_\allowbreak{}natural & decimal & 16 & 0 & 0.669551 \\
 & tot\_\allowbreak{}pob\_\allowbreak{}internet & decimal & 16 & 0 & 0.60124 \\
 & tot\_\allowbreak{}pob\_\allowbreak{}ipm & decimal & 16 & 0 & 0.101172 \\
 & tot\_\allowbreak{}pob\_\allowbreak{}ipm\_\allowbreak{}hogares & decimal & 16 & 0 & 0.0786242 \\
 & tot\_\allowbreak{}pob\_\allowbreak{}nbi & decimal & 16 & 0 & 0.11355 \\
 & tot\_\allowbreak{}pob\_\allowbreak{}recoleccion\_\allowbreak{}basuras & decimal & 16 & 0 & 0.944698 \\
 & tot\_\allowbreak{}td & decimal & 11 & 37.5 & 0.0740896 \\
 & tot\_\allowbreak{}td\_\allowbreak{}15\_\allowbreak{}28 & decimal & 11 & 37.5 & 0.105271 \\
 & tot\_\allowbreak{}to & decimal & 11 & 37.5 & 0.510432 \\
 & urb\_\allowbreak{}emp\_\allowbreak{}informal & decimal & 11 & 37.5 & 0.44519 \\
 & urb\_\allowbreak{}gini\_\allowbreak{}hog & decimal & 16 & 6.2 & 0.436235 \\
 & urb\_\allowbreak{}gini\_\allowbreak{}lab & decimal & 16 & 6.2 & 0.327307 \\
 & urb\_\allowbreak{}nini & decimal & 16 & 0 & 0.234975 \\
 & urb\_\allowbreak{}pob\_\allowbreak{}acueducto & decimal & 16 & 0 & 0.975146 \\
 & urb\_\allowbreak{}pob\_\allowbreak{}alcantarillado & decimal & 16 & 0 & 0.960052 \\
 & urb\_\allowbreak{}pob\_\allowbreak{}energia & decimal & 16 & 0 & 0.995155 \\
 & urb\_\allowbreak{}pob\_\allowbreak{}gas\_\allowbreak{}natural & decimal & 16 & 0 & 0.776224 \\
 & urb\_\allowbreak{}pob\_\allowbreak{}internet & decimal & 16 & 0 & 0.656756 \\
 & urb\_\allowbreak{}pob\_\allowbreak{}ipm & decimal & 16 & 0 & 0.0768115 \\
 & urb\_\allowbreak{}pob\_\allowbreak{}ipm\_\allowbreak{}hogares & decimal & 16 & 0 & 0.0572301 \\
 & urb\_\allowbreak{}pob\_\allowbreak{}nbi & decimal & 16 & 0 & 0.100781 \\
 & urb\_\allowbreak{}pob\_\allowbreak{}recoleccion\_\allowbreak{}basuras & decimal & 16 & 0 & 0.976035 \\
 & urb\_\allowbreak{}td & decimal & 11 & 37.5 & 0.0767988 \\
 & urb\_\allowbreak{}td\_\allowbreak{}15\_\allowbreak{}28 & decimal & 11 & 37.5 & 0.111099 \\
 & urb\_\allowbreak{}to & decimal & 11 & 37.5 & 0.515796 \\
ECV\_\allowbreak{}raw & Municipio & decimal & 125 & 0 &    5001 \\
 & NomMunicipio & texto & 125 & 0 & Medellín \\
 & Indicador & texto & 794 & 0 & SPAL \\
 & NomIndicador & texto & 791 & 0 & Porcentaje de viviendas c... \\
 & Valor\_\allowbreak{}tot & decimal & 74503 & 1.7 & 97.9849 \\
 & SE\_\allowbreak{}tot & decimal & 72454 & 1.7 & 0.155248 \\
 & IC25\_\allowbreak{}tot & decimal & 56515 & 1.7 & 97.6806 \\
 & IC975\_\allowbreak{}tot & decimal & 75082 & 1.7 & 98.2892 \\
 & CV\_\allowbreak{}tot & decimal & 72266 & 19.5 & 0.15844 \\
 & Clase\_\allowbreak{}tot & texto & 1 & 0 & Total \\
 & Valor\_\allowbreak{}urb & decimal & 56894 & 2.6 & 99.0921 \\
 & SE\_\allowbreak{}urb & decimal & 61007 & 2.6 & 0.114786 \\
 & IC25\_\allowbreak{}urb & decimal & 46206 & 2.6 & 98.8671 \\
 & IC975\_\allowbreak{}urb & decimal & 64149 & 2.6 & 99.3171 \\
 & CV\_\allowbreak{}urb & decimal & 61866 & 25.7 & 0.115838 \\
 & Clase\_\allowbreak{}urb & texto & 2 & 0 & Urbano \\
 & Valor\_\allowbreak{}rur & decimal & 52171 & 3.5 &  90.184 \\
 & SE\_\allowbreak{}rur & decimal & 55856 & 3.5 & 0.951883 \\
 & IC25\_\allowbreak{}rur & decimal & 41135 & 3.5 & 88.3184 \\
 & IC975\_\allowbreak{}rur & decimal & 58878 & 3.5 & 92.0497 \\
 & CV\_\allowbreak{}rur & decimal & 56258 & 29.8 & 1.05549 \\
 & Clase\_\allowbreak{}rur & texto & 1 & 0 & Rural \\
 & Tipo & texto & 3 & 0 & Viviendas \\
 & GrupoPoblacional & texto & 3 & 0 & General \\
 & Y & decimal & 1 & 100 &      NA \\
 & Z & texto & 124 & 0 & 05001 \\
estructura\_\allowbreak{}productiva & ind\_\allowbreak{}mpio & texto & 135 & 0 & 05 \\
 & creacion\_\allowbreak{}emp & decimal & 95 & 0.7 &   23537 \\
 & dens\_\allowbreak{}emp & decimal & 135 & 0 & 33.4093 \\
 & fin\_\allowbreak{}banc\_\allowbreak{}pc & decimal & 126 & 7.4 & 6.75606 \\
 & fin\_\allowbreak{}microcr\_\allowbreak{}pc & decimal & 116 & 7.4 & 3.31638 \\
 & tnat\_\allowbreak{}emp & decimal & 126 & 7.4 &  17.645 \\
 & va\_\allowbreak{}pc & decimal & 126 & 7.4 & 21625.3 \\
 & va\_\allowbreak{}primario & decimal & 126 & 7.4 & 0.195439 \\
imca\_\allowbreak{}2022 & ind\_\allowbreak{}mpio & texto & 134 & 0 & 05001 \\
 & imca\_\allowbreak{}adop\_\allowbreak{}tic & decimal & 134 & 0 & 50.4179 \\
 & imca\_\allowbreak{}capacidades & decimal & 134 & 0 & 56.8596 \\
 & imca\_\allowbreak{}dinam\_\allowbreak{}neg & decimal & 134 & 0 & 86.1516 \\
 & imca\_\allowbreak{}infraestructura & decimal & 134 & 0 & 99.2887 \\
 & imca\_\allowbreak{}innovacion & decimal & 16 & 0 & 20.3998 \\
 & imca\_\allowbreak{}instituciones & decimal & 134 & 0 & 71.0171 \\
 & imca\_\allowbreak{}merc\_\allowbreak{}bienes & decimal & 134 & 0 & 72.8602 \\
 & imca\_\allowbreak{}merc\_\allowbreak{}laboral & decimal & 134 & 0 &  52.872 \\
 & imca\_\allowbreak{}salud & decimal & 134 & 0 & 57.2265 \\
 & imca\_\allowbreak{}sist\_\allowbreak{}financiero & decimal & 134 & 0 & 51.8916 \\
 & imca\_\allowbreak{}tam\_\allowbreak{}mercado & decimal & 134 & 0 & 84.5544 \\
 & imca\_\allowbreak{}total & decimal & 134 & 0 & 63.9581 \\
 & Año & decimal & 1 & 0 &    2022 \\
infraestructura\_\allowbreak{}internet\_\allowbreak{}2025 & anno & decimal & 1 & 0 &    2025 \\
 & trimestre & decimal & 4 & 0 &       1 \\
 & id\_\allowbreak{}empresa & decimal & 286 & 0 & 8.00137e+08 \\
 & empresa & texto & 286 & 0 & CENTURYLINK COLOMBIA S.A. \\
 & ind\_\allowbreak{}mpio & decimal & 125 & 0 &    5093 \\
 & municipio & texto & 125 & 0 & BETULIA \\
 & id\_\allowbreak{}departamento & decimal & 1 & 0 &       5 \\
 & departamento & texto & 1 & 0 & ANTIOQUIA \\
 & id\_\allowbreak{}segmento & decimal & 10 & 0 &     107 \\
 & segmento & texto & 10 & 0 & Corporativo \\
 & id\_\allowbreak{}servicio\_\allowbreak{}paquete & decimal & 7 & 0 &       1 \\
 & servicio\_\allowbreak{}paquete & texto & 7 & 0 & Internet fijo \\
 & velocidad\_\allowbreak{}efectiva\_\allowbreak{}downstream & texto & 490 & 0 & 200 \\
 & velocidad\_\allowbreak{}efectiva\_\allowbreak{}upstream & texto & 589 & 0 & 200 \\
 & id\_\allowbreak{}tecnologia & decimal & 15 & 0 &     107 \\
 & tecnologia & texto & 15 & 0 & Otras tecnologías de fibr... \\
 & id\_\allowbreak{}estado & decimal & 2 & 0 &       1 \\
 & estado & texto & 2 & 0 & Activo en funcionamiento \\
 & cantidad\_\allowbreak{}lineas\_\allowbreak{}accesos & decimal & 2017 & 0 &       1 \\
 & valor\_\allowbreak{}facturado\_\allowbreak{}o\_\allowbreak{}cobrado & decimal & 76906 & 0 & 1.20535e+07 \\
 & otros\_\allowbreak{}valores\_\allowbreak{}facturados & decimal & 40497 & 0 &       0 \\
 & valor\_\allowbreak{}total\_\allowbreak{}plan\_\allowbreak{}tarifario & decimal & 1 & 0 &       0 \\
 & period & decimal & 4 & 0 &     260 \\
percepcion\_\allowbreak{}seguridad & ESTUDIO & texto & 8 & 0 & 079518000010 \\
 & FECHAINI & texto & 256 & 0 & 29/10/2018 \\
 & B\_\allowbreak{}SUBREGION & decimal & 9 & 0 &       9 \\
 & P1A & decimal & 4 & 0 &       2 \\
 & P2 & decimal & 4 & 0 &       3 \\
 & P4A & decimal & 4 & 0 &       2 \\
 & P6 & decimal & 4 & 0 &       3 \\
 & FACTOR\_\allowbreak{}PONDERACION & decimal & 2972 & 0 & 4.70912 \\
poblacion\_\allowbreak{}municipal\_\allowbreak{}total\_\allowbreak{}2025 & ind\_\allowbreak{}mpio & decimal & 1123 & 0 &    5001 \\
 & nvl\_\allowbreak{}label & texto & 1040 & 0 & Medellín \\
 & P\_\allowbreak{}T & decimal & 2 & 0.1 &     100 \\
 & P\_\allowbreak{}H & decimal & 1122 & 0.1 & 47.2697 \\
 & P\_\allowbreak{}M & decimal & 1122 & 0.1 & 52.7303 \\
 & P\_\allowbreak{}inf\_\allowbreak{}T\_\allowbreak{}0\_\allowbreak{}11 & decimal & 1123 & 0.1 & 13.1976 \\
 & P\_\allowbreak{}j\_\allowbreak{}T\_\allowbreak{}12\_\allowbreak{}28 & decimal & 1121 & 0.1 & 25.1281 \\
 & P\_\allowbreak{}ad\_\allowbreak{}T\_\allowbreak{}29\_\allowbreak{}59 & decimal & 1123 & 0.1 & 43.5166 \\
 & P\_\allowbreak{}v\_\allowbreak{}T\_\allowbreak{}60\_\allowbreak{}85 & decimal & 1123 & 0.1 & 18.1577 \\
 & P\_\allowbreak{}inf\_\allowbreak{}H\_\allowbreak{}0\_\allowbreak{}11 & decimal & 1122 & 0.1 & 6.73734 \\
 & P\_\allowbreak{}j\_\allowbreak{}H\_\allowbreak{}12\_\allowbreak{}28 & decimal & 1121 & 0.1 & 12.6658 \\
 & P\_\allowbreak{}ad\_\allowbreak{}H\_\allowbreak{}29\_\allowbreak{}59 & decimal & 1123 & 0.1 &  20.647 \\
 & P\_\allowbreak{}v\_\allowbreak{}H\_\allowbreak{}60\_\allowbreak{}85 & decimal & 1123 & 0.1 & 7.21963 \\
 & P\_\allowbreak{}inf\_\allowbreak{}M\_\allowbreak{}0\_\allowbreak{}11 & decimal & 1123 & 0.1 & 6.46024 \\
 & P\_\allowbreak{}j\_\allowbreak{}M\_\allowbreak{}12\_\allowbreak{}28 & decimal & 1120 & 0.1 & 12.4624 \\
 & P\_\allowbreak{}ad\_\allowbreak{}M\_\allowbreak{}29\_\allowbreak{}59 & decimal & 1122 & 0.1 & 22.8697 \\
 & P\_\allowbreak{}v\_\allowbreak{}M\_\allowbreak{}60\_\allowbreak{}85 & decimal & 1123 & 0.1 &  10.938 \\
 & IDE\_\allowbreak{}T & decimal & 1123 & 0.1 & 54.2868 \\
 & IDE\_\allowbreak{}H & decimal & 1122 & 0.1 & 50.7293 \\
 & IDE\_\allowbreak{}M & decimal & 1120 & 0.1 & 57.6218 \\
 & I\_\allowbreak{}enve\_\allowbreak{}T & decimal & 1123 & 0.1 & 36.1417 \\
 & I\_\allowbreak{}enve\_\allowbreak{}H & decimal & 1120 & 0.1 & 30.8864 \\
 & I\_\allowbreak{}enve\_\allowbreak{}M & decimal & 1122 & 0.1 &  40.479 \\
 & P\_\allowbreak{}T\_\allowbreak{}r & decimal & 2 & 0.1 &     100 \\
 & P\_\allowbreak{}H\_\allowbreak{}r & decimal & 1121 & 0.1 & 50.7871 \\
 & P\_\allowbreak{}M\_\allowbreak{}r & decimal & 1121 & 0.1 & 49.2129 \\
 & P\_\allowbreak{}inf\_\allowbreak{}T\_\allowbreak{}0\_\allowbreak{}11\_\allowbreak{}r & decimal & 1121 & 0.1 & 15.1644 \\
 & P\_\allowbreak{}j\_\allowbreak{}T\_\allowbreak{}12\_\allowbreak{}28\_\allowbreak{}r & decimal & 1122 & 0.1 & 24.5479 \\
 & P\_\allowbreak{}ad\_\allowbreak{}T\_\allowbreak{}29\_\allowbreak{}59\_\allowbreak{}r & decimal & 1122 & 0.1 & 44.0506 \\
 & P\_\allowbreak{}v\_\allowbreak{}T\_\allowbreak{}60\_\allowbreak{}85\_\allowbreak{}r & decimal & 1122 & 0.1 & 16.2372 \\
 & P\_\allowbreak{}inf\_\allowbreak{}H\_\allowbreak{}0\_\allowbreak{}11\_\allowbreak{}r & decimal & 1121 & 0.1 & 7.70243 \\
 & P\_\allowbreak{}j\_\allowbreak{}H\_\allowbreak{}12\_\allowbreak{}28\_\allowbreak{}r & decimal & 1123 & 0.1 & 12.6901 \\
 & P\_\allowbreak{}ad\_\allowbreak{}H\_\allowbreak{}29\_\allowbreak{}59\_\allowbreak{}r & decimal & 1121 & 0.1 & 22.3777 \\
 & P\_\allowbreak{}v\_\allowbreak{}H\_\allowbreak{}60\_\allowbreak{}85\_\allowbreak{}r & decimal & 1122 & 0.1 & 8.01685 \\
 & P\_\allowbreak{}inf\_\allowbreak{}M\_\allowbreak{}0\_\allowbreak{}11\_\allowbreak{}r & decimal & 1121 & 0.1 & 7.46198 \\
 & P\_\allowbreak{}j\_\allowbreak{}M\_\allowbreak{}12\_\allowbreak{}28\_\allowbreak{}r & decimal & 1119 & 0.1 & 11.8578 \\
 & P\_\allowbreak{}ad\_\allowbreak{}M\_\allowbreak{}29\_\allowbreak{}59\_\allowbreak{}r & decimal & 1123 & 0.1 & 21.6728 \\
 & P\_\allowbreak{}v\_\allowbreak{}M\_\allowbreak{}60\_\allowbreak{}85\_\allowbreak{}r & decimal & 1123 & 0.1 &  8.2203 \\
 & IDE\_\allowbreak{}T\_\allowbreak{}r & decimal & 1123 & 0.1 & 55.1807 \\
 & IDE\_\allowbreak{}H\_\allowbreak{}r & decimal & 1120 & 0.1 & 54.5432 \\
 & IDE\_\allowbreak{}M\_\allowbreak{}r & decimal & 1121 & 0.1 & 55.8441 \\
 & I\_\allowbreak{}enve\_\allowbreak{}T\_\allowbreak{}r & decimal & 1122 & 0.1 & 30.6883 \\
 & I\_\allowbreak{}enve\_\allowbreak{}H\_\allowbreak{}r & decimal & 1118 & 0.1 &  29.523 \\
 & I\_\allowbreak{}enve\_\allowbreak{}M\_\allowbreak{}r & decimal & 1119 & 0.1 & 31.8727 \\
 & P\_\allowbreak{}T\_\allowbreak{}c & decimal & 2 & 1.7 &     100 \\
 & P\_\allowbreak{}H\_\allowbreak{}c & decimal & 1101 & 1.7 & 47.2007 \\
 & P\_\allowbreak{}M\_\allowbreak{}c & decimal & 1101 & 1.7 & 52.7993 \\
 & P\_\allowbreak{}inf\_\allowbreak{}T\_\allowbreak{}0\_\allowbreak{}11\_\allowbreak{}c & decimal & 1099 & 1.7 &  13.159 \\
 & P\_\allowbreak{}j\_\allowbreak{}T\_\allowbreak{}12\_\allowbreak{}28\_\allowbreak{}c & decimal & 1102 & 1.7 & 25.1395 \\
 & P\_\allowbreak{}ad\_\allowbreak{}T\_\allowbreak{}29\_\allowbreak{}59\_\allowbreak{}c & decimal & 1102 & 1.7 & 43.5062 \\
 & P\_\allowbreak{}v\_\allowbreak{}T\_\allowbreak{}60\_\allowbreak{}85\_\allowbreak{}c & decimal & 1105 & 1.7 & 18.1954 \\
 & P\_\allowbreak{}inf\_\allowbreak{}H\_\allowbreak{}0\_\allowbreak{}11\_\allowbreak{}c & decimal & 1102 & 1.7 &  6.7184 \\
 & P\_\allowbreak{}j\_\allowbreak{}H\_\allowbreak{}12\_\allowbreak{}28\_\allowbreak{}c & decimal & 1102 & 1.7 & 12.6653 \\
 & P\_\allowbreak{}ad\_\allowbreak{}H\_\allowbreak{}29\_\allowbreak{}59\_\allowbreak{}c & decimal & 1101 & 1.7 &  20.613 \\
 & P\_\allowbreak{}v\_\allowbreak{}H\_\allowbreak{}60\_\allowbreak{}85\_\allowbreak{}c & decimal & 1100 & 1.7 & 7.20399 \\
 & P\_\allowbreak{}inf\_\allowbreak{}M\_\allowbreak{}0\_\allowbreak{}11\_\allowbreak{}c & decimal & 1099 & 1.7 & 6.44058 \\
 & P\_\allowbreak{}j\_\allowbreak{}M\_\allowbreak{}12\_\allowbreak{}28\_\allowbreak{}c & decimal & 1099 & 1.7 & 12.4742 \\
 & P\_\allowbreak{}ad\_\allowbreak{}M\_\allowbreak{}29\_\allowbreak{}59\_\allowbreak{}c & decimal & 1102 & 1.7 & 22.8932 \\
 & P\_\allowbreak{}v\_\allowbreak{}M\_\allowbreak{}60\_\allowbreak{}85\_\allowbreak{}c & decimal & 1104 & 1.7 & 10.9914 \\
 & IDE\_\allowbreak{}T\_\allowbreak{}c & decimal & 1102 & 1.7 & 54.2694 \\
 & IDE\_\allowbreak{}H\_\allowbreak{}c & decimal & 1098 & 1.7 & 50.6508 \\
 & IDE\_\allowbreak{}M\_\allowbreak{}c & decimal & 1099 & 1.7 & 57.6547 \\
 & I\_\allowbreak{}enve\_\allowbreak{}c\_\allowbreak{}T & decimal & 1103 & 1.7 & 36.2499 \\
 & I\_\allowbreak{}enve\_\allowbreak{}H\_\allowbreak{}c & decimal & 1095 & 1.7 & 30.9166 \\
 & I\_\allowbreak{}enve\_\allowbreak{}M\_\allowbreak{}c & decimal & 1095 & 1.7 & 40.6332 \\
 & PR\_\allowbreak{}T & decimal & 1106 & 0.1 & 1.92458 \\
 & PR\_\allowbreak{}H & decimal & 1105 & 0.1 & 2.06779 \\
 & PR\_\allowbreak{}M & decimal & 1106 & 0.1 &  1.7962 \\
 & PR\_\allowbreak{}inf\_\allowbreak{}T\_\allowbreak{}0\_\allowbreak{}11 & decimal & 1105 & 0.1 &  2.2114 \\
 & PR\_\allowbreak{}j\_\allowbreak{}T\_\allowbreak{}12\_\allowbreak{}28 & decimal & 1106 & 0.1 & 1.88014 \\
 & PR\_\allowbreak{}ad\_\allowbreak{}T\_\allowbreak{}29\_\allowbreak{}59 & decimal & 1105 & 0.1 & 1.94819 \\
 & PR\_\allowbreak{}v\_\allowbreak{}T\_\allowbreak{}60\_\allowbreak{}85 & decimal & 1104 & 0.1 & 1.72102 \\
 & PR\_\allowbreak{}inf\_\allowbreak{}H\_\allowbreak{}0\_\allowbreak{}11 & decimal & 1102 & 0.1 & 2.20027 \\
 & PR\_\allowbreak{}j\_\allowbreak{}H\_\allowbreak{}12\_\allowbreak{}28 & decimal & 1105 & 0.1 & 1.92828 \\
 & PR\_\allowbreak{}ad\_\allowbreak{}H\_\allowbreak{}29\_\allowbreak{}59 & decimal & 1104 & 0.1 & 2.08591 \\
 & PR\_\allowbreak{}v\_\allowbreak{}H\_\allowbreak{}60\_\allowbreak{}85 & decimal & 1102 & 0.1 &  2.1371 \\
 & PR\_\allowbreak{}inf\_\allowbreak{}M\_\allowbreak{}0\_\allowbreak{}11 & decimal & 1100 & 0.1 & 2.22301 \\
 & PR\_\allowbreak{}j\_\allowbreak{}M\_\allowbreak{}12\_\allowbreak{}28 & decimal & 1103 & 0.1 & 1.83122 \\
 & PR\_\allowbreak{}ad\_\allowbreak{}M\_\allowbreak{}29\_\allowbreak{}59 & decimal & 1100 & 0.1 & 1.82386 \\
 & PR\_\allowbreak{}v\_\allowbreak{}M\_\allowbreak{}60\_\allowbreak{}85 & decimal & 1101 & 0.1 & 1.44639 \\
 & PU\_\allowbreak{}T & decimal & 1106 & 0.1 & 98.0754 \\
 & PU\_\allowbreak{}H & decimal & 1105 & 0.1 & 97.9322 \\
 & PU\_\allowbreak{}M & decimal & 1106 & 0.1 & 98.2038 \\
 & PU\_\allowbreak{}inf\_\allowbreak{}T\_\allowbreak{}0\_\allowbreak{}11 & decimal & 1105 & 0.1 & 97.7886 \\
 & PU\_\allowbreak{}j\_\allowbreak{}T\_\allowbreak{}12\_\allowbreak{}28 & decimal & 1106 & 0.1 & 98.1199 \\
 & PU\_\allowbreak{}ad\_\allowbreak{}T\_\allowbreak{}29\_\allowbreak{}59 & decimal & 1105 & 0.1 & 98.0518 \\
 & PU\_\allowbreak{}v\_\allowbreak{}T\_\allowbreak{}60\_\allowbreak{}85 & decimal & 1104 & 0.1 &  98.279 \\
 & PU\_\allowbreak{}inf\_\allowbreak{}H\_\allowbreak{}0\_\allowbreak{}11 & decimal & 1102 & 0.1 & 97.7997 \\
 & PU\_\allowbreak{}j\_\allowbreak{}H\_\allowbreak{}12\_\allowbreak{}28 & decimal & 1105 & 0.1 & 98.0717 \\
 & PU\_\allowbreak{}ad\_\allowbreak{}H\_\allowbreak{}29\_\allowbreak{}59 & decimal & 1104 & 0.1 & 97.9141 \\
 & PU\_\allowbreak{}v\_\allowbreak{}H\_\allowbreak{}60\_\allowbreak{}85 & decimal & 1102 & 0.1 & 97.8629 \\
 & PU\_\allowbreak{}inf\_\allowbreak{}M\_\allowbreak{}0\_\allowbreak{}11 & decimal & 1100 & 0.1 &  97.777 \\
 & PU\_\allowbreak{}j\_\allowbreak{}M\_\allowbreak{}12\_\allowbreak{}28 & decimal & 1103 & 0.1 & 98.1688 \\
 & PU\_\allowbreak{}ad\_\allowbreak{}M\_\allowbreak{}29\_\allowbreak{}59 & decimal & 1100 & 0.1 & 98.1761 \\
 & PU\_\allowbreak{}v\_\allowbreak{}M\_\allowbreak{}60\_\allowbreak{}85 & decimal & 1101 & 0.1 & 98.5536 \\
poblacion\_\allowbreak{}total\_\allowbreak{}2025 & cod\_\allowbreak{}dpto & texto & 1 & 0 & 05 \\
 & departamento & texto & 1 & 0 & Antioquia \\
 & ind\_\allowbreak{}mpio & decimal & 125 & 0 &    5001 \\
 & nvl\_\allowbreak{}label & texto & 125 & 0 & Medellín \\
 & año & decimal & 1 & 0 &    2025 \\
 & area\_\allowbreak{}geo & texto & 3 & 0 & Cabecera Municipal \\
 & Total & decimal & 374 & 0 & 2.47968e+06 \\
 & Hombres & decimal & 369 & 0 & 1.17043e+06 \\
 & Mujeres & decimal & 369 & 0 & 1.30926e+06 \\
seguridad\_\allowbreak{}policia & ind\_\allowbreak{}mpio & decimal & 125 & 0 &    5001 \\
 & delito & texto & 4 & 0 & homicidios \\
 & anio & entero & 2 & 0 &    2025 \\
 & cod\_\allowbreak{}mpio & texto & 125 & 0 & 05001 \\
 & nvl\_\allowbreak{}label & texto & 125 & 0 & Medellín \\
 & total\_\allowbreak{}historico & decimal & 254 & 0 &    1783 \\
 & total\_\allowbreak{}2025 & decimal & 111 & 50 &     333 \\
 & arma\_\allowbreak{}fuego\_\allowbreak{}2025 & decimal & 45 & 50 &     155 \\
 & arma\_\allowbreak{}blanca\_\allowbreak{}2025 & decimal & 31 & 50 &     127 \\
 & contundente\_\allowbreak{}2025 & decimal & 44 & 50 &      51 \\
 & dron\_\allowbreak{}explosivo\_\allowbreak{}2025 & decimal & 2 & 87.5 &       0 \\
 & otros\_\allowbreak{}2025 & decimal & 16 & 50 &       0 \\
 & sin\_\allowbreak{}armas\_\allowbreak{}2025 & decimal & 90 & 62.5 &    8081 \\
 & total\_\allowbreak{}2023 & decimal & 110 & 50 &     359 \\
 & arma\_\allowbreak{}fuego\_\allowbreak{}2023 & decimal & 49 & 50 &     152 \\
 & arma\_\allowbreak{}blanca\_\allowbreak{}2023 & decimal & 33 & 50 &     154 \\
 & contundente\_\allowbreak{}2023 & decimal & 43 & 50 &      53 \\
 & dron\_\allowbreak{}explosivo\_\allowbreak{}2023 & decimal & 2 & 87.5 &       0 \\
 & otros\_\allowbreak{}2023 & decimal & 20 & 50 &       0 \\
 & sin\_\allowbreak{}armas\_\allowbreak{}2023 & decimal & 90 & 62.5 &    7955 \\
subreg\_\allowbreak{}completo & ind\_\allowbreak{}mpio & entero & 125 & 0 &    5001 \\
 & subregion\_\allowbreak{}full & texto & 9 & 0 & VALLE DE ABURRA \\
uso\_\allowbreak{}del\_\allowbreak{}suelo\_\allowbreak{}pota & codmpio & texto & 88 & 0 & 05002 \\
 & municipio & texto & 88 & 0 & abejorral \\
 & grupo & texto & 8 & 0 & GRUPO 3 \\
 & n\_\allowbreak{}predios\_\allowbreak{}rural & entero & 86 & 0 &    6996 \\
 & pct\_\allowbreak{}uso\_\allowbreak{}adecuado\_\allowbreak{}rural & decimal & 88 & 0 & 26.3032 \\
 & pct\_\allowbreak{}sobreutilizacion\_\allowbreak{}rural & decimal & 88 & 0 & 68.1716 \\
 & pct\_\allowbreak{}subutilizacion\_\allowbreak{}rural & decimal & 88 & 0 & 5.02341 \\
 & pct\_\allowbreak{}sin\_\allowbreak{}conflicto\_\allowbreak{}rural & decimal & 88 & 0 & 0.501865 \\
 & chk\_\allowbreak{}pct\_\allowbreak{}sum & decimal & 37 & 0 &     100 \\
 & n\_\allowbreak{}predios\_\allowbreak{}uso\_\allowbreak{}adec\_\allowbreak{}lt60 & entero & 88 & 0 &    6160 \\
 & n\_\allowbreak{}predios\_\allowbreak{}uso\_\allowbreak{}adec\_\allowbreak{}60\_\allowbreak{}80 & entero & 76 & 0 &     300 \\
 & n\_\allowbreak{}predios\_\allowbreak{}uso\_\allowbreak{}adec\_\allowbreak{}gt80 & entero & 83 & 0 &     536 \\
 & n\_\allowbreak{}predios\_\allowbreak{}uso\_\allowbreak{}adec\_\allowbreak{}100 & entero & 82 & 0 &     173 \\
 & tiene\_\allowbreak{}urbano & decimal & 1 & 0 &       1 \\
 & n\_\allowbreak{}predios\_\allowbreak{}urbano & entero & 67 & 0 &    3045 \\
 & n\_\allowbreak{}predios\_\allowbreak{}favorables\_\allowbreak{}proteccion & entero & 39 & 0 &     252 \\
 & n\_\allowbreak{}predios\_\allowbreak{}no\_\allowbreak{}favorables & entero & 80 & 0 &    2793 \\
\end{longtable}

}

\end{document}
